\chapter{回溯、枚举和搜索空间剪枝}

回溯题表面上是在“试所有可能”，本质上是在一棵决策树上搜索。每一层做一个选择，每条从根到叶子的路径代表一个候选答案。真正要学会的不是某个递归模板，而是四个问题：当前层决定什么，已经选择了什么，下一步还能选什么，什么时候可以停止或剪枝。只要这四件事说不清，代码再短也只是碰巧通过样例。

暴力枚举和回溯的边界并不绝对。最朴素的暴力会把所有组合全部生成出来，再逐个检查是否合法；回溯则把“合法性判断”提前到生成过程中。一旦发现当前前缀已经不可能形成答案，就不再向下扩展。这个动作叫剪枝。剪枝不是为了让代码看起来高级，而是减少决策树里根本不需要访问的分支。

\cpp 面试题常用递归表达回溯，因为调用栈天然保存了“走到第几层、下一步尝试哪个分支、回退后恢复什么状态”。但递归深度不受控时会带来工程风险。本章讲递归心智模型，代码尽量把栈帧显式写出来，让搜索过程、状态恢复和边界条件更可见。这样写会比递归略长，但它能训练你看清楚回溯真正维护的状态。

本章先讲搜索树、路径状态、候选集合、剪枝和去重，再用排列、子集、组合、括号、棋盘和字符串切分题展示这些概念如何落地。回溯题卡现在单独放在本章末尾；位运算和状态压缩作为枚举状态的一种表示，也放在回溯之后，而不是混在哈希或 DP 章节里。

\section{搜索树、约束和输出规模}

回溯题要先画搜索树。树的第几层代表什么，取决于题目对象。全排列中，第 \code{depth} 层决定排列的第 \code{depth} 个位置；子集中，第 \code{index} 层决定当前元素选或不选；组合题中，一层选择下一个加入组合的候选；切分题中，一层选择下一个切分终点；网格路径中，一层选择下一步方向。层含义不同，代码里的 \code{start}、\code{used}、\code{index} 和访问标记就不同，不能统一套一个函数形状。

路径状态表示从根到当前节点已经做出的选择。它可以是一个数组、一段字符串、若干个切分片段、棋盘上已经放置的皇后位置，或者一个网格路径访问集合。候选集合表示当前层还能选什么。排列题的候选是所有未使用下标；组合题的候选是从 \code{start} 往后的下标；括号生成的候选是左括号或右括号，但受剩余数量限制；N 皇后的候选是当前行的各列，但受列和两条对角线限制。把路径和候选集合分清，回溯代码才有明确职责。

剪枝必须是必要条件或支配关系，不能是直觉。先看组合总和：候选 \code{[2,3,6,7]} 已排序、剩余目标是 5 时，试到 6 已经超出，后面的 7 更不可能使用，所以可以整段停止；如果候选里允许负数，这个推理就失效，因为后面可能用负数把总和拉回来。括号生成中，前缀 \code{"())"} 已经右括号更多，后面再补左括号也无法把这个非法前缀变成合法括号串；复原 IP 地址中，剩 2 段却剩 7 个字符，也不可能每段最多 3 个字符装下。N 皇后中，同列或同对角线冲突一定非法。没有证明的剪枝会制造隐藏漏解，尤其是在输入有负数、重复值或目标条件变化时。

去重也要分层次。路径内去重解决“同一个元素不能在一条路径里重复使用”，常用 \code{used} 或访问标记；同层去重解决“相同值作为本层不同分支会生成重复答案”，常用排序后跳过相邻重复值。全排列 II 同时有路径内去重和同层去重；组合总和 II 通过 \code{start} 保证下标不重复，再通过同层跳过重复值保证答案不重复；子集 II 的去重只在同一层发生。把这些条件混在一起，会出现要么漏掉合法重复使用，要么产生重复答案。

输出规模决定复杂度下限。全排列有 \complexity{n!} 个答案，子集有 \complexity{2^n} 个答案，括号生成有卡特兰数个答案，N 皇后答案数量也随 n 快速增长。如果题目要求输出所有方案，算法时间不可能低于输出总大小。复杂度分析要把“搜索访问的状态数”和“复制答案的成本”分开：生成一个长度为 n 的排列，光复制路径就是 \complexity{O(n)}。这不是实现低效，而是输出本身需要这么多信息。

递归和迭代只是保存栈帧的两种方式。递归函数参数通常包括层号、路径、候选起点、剩余目标和访问标记；显式栈版本则把这些字段放进结构体。递归写法短，适合表达；显式栈写法长，但能避免调用栈风险，也更容易看清状态复制成本。本书代码偏向显式栈，并不否认递归思维；恰恰相反，只有先理解递归栈帧保存了什么，才能把它可靠地改写成迭代。

\begin{principlebox}{回溯题的建模清单}
先写搜索树层含义，再写路径状态和候选集合；然后写终止条件，说明何时把路径加入答案；接着写剪枝条件，并证明被剪掉的分支不可能产生答案；最后写去重条件，区分路径内去重和同层去重。若答案数量本身是指数级，复杂度表达必须包含输出规模。
\end{principlebox}

全排列是回溯的第一道标准题。题目给定互不相同的数字，要求返回所有排列。暴力想法是枚举长度为 \code{n} 的所有序列，再检查是否每个数字恰好使用一次。若把每个位置都看成有 \code{n} 种选择，总分支数是 \complexity{n^n}，绝大多数序列会因为重复使用元素而被丢弃。优化方法是在生成过程中维护 \code{used} 数组，保证同一个下标不会被选两次。这样搜索树只保留真正的排列分支，叶子数变成 \complexity{n!}。

下面的实现用显式栈模拟递归。每个栈帧保存当前层接下来要尝试的候选下标 \code{next_index}。进入下一层时，把选择加入 \code{path} 并标记 \code{used}；一层候选全部尝试完时，弹出栈帧，并撤销上一层做出的选择。

\begin{lstlisting}[language=C++,caption={LeetCode 46 全排列：显式栈模拟回溯}]
std::vector<std::vector<int>> permute(const std::vector<int>& nums)
{
    struct Frame {
        int next_index = 0;
    };

    const int n = static_cast<int>(nums.size());
    std::vector<std::vector<int>> answer;
    std::vector<int> path;
    std::vector<char> used(n, false);
    std::vector<Frame> stack{{0}};

    while (!stack.empty()) {
        if (static_cast<int>(path.size()) == n) {
            answer.push_back(path);
            stack.pop_back();
            if (!path.empty()) {
                const int last_value = path.back();
                path.pop_back();
                const auto found = std::find(nums.begin(), nums.end(), last_value);
                used[static_cast<int>(found - nums.begin())] = false;
            }
            continue;
        }

        Frame& frame = stack.back();
        while (frame.next_index < n && used[frame.next_index]) {
            ++frame.next_index;
        }
        if (frame.next_index == n) {
            stack.pop_back();
            if (!path.empty()) {
                const int last_value = path.back();
                path.pop_back();
                const auto found = std::find(nums.begin(), nums.end(), last_value);
                used[static_cast<int>(found - nums.begin())] = false;
            }
            continue;
        }

        const int chosen_index = frame.next_index;
        ++frame.next_index;
        used[chosen_index] = true;
        path.push_back(nums[chosen_index]);
        stack.push_back(Frame{0});
    }

    return answer;
}
\end{lstlisting}

这段代码为了突出栈帧，撤销选择时用值反查下标；在生产代码里可以把“本层选择的下标”也放进下一层栈帧，避免反查。它的时间复杂度必须按输出规模理解：一共有 \complexity{n!} 个排列，每个排列复制长度为 \code{n} 的路径，所以输出成本是 \complexity{O(n \cdot n!)}。额外空间是路径、标记和栈，都是 \complexity{O(n)}，不计答案。复盘这题时要能说清楚：\code{used} 是路径合法性的约束，\code{path.size() == n} 是叶子条件，回退必须撤销最后一次选择。

如果输入含重复数字，全排列 II 的难点就变成去重。很多人会先生成所有排列，再放进 \code{set} 去重；这当然能做，但浪费大量重复分支。更好的方法是先排序，然后在同一层跳过“相同数值且前一个相同值还没有被本路径使用”的候选。这个条件的含义是：同一层里，相同值只允许最左边那个代表这一类选择，避免生成形如交换两个相同值位置的重复排列。

\begin{lstlisting}[language=C++,caption={LeetCode 47 全排列 II：排序后按层去重}]
std::vector<std::vector<int>> permute_unique(std::vector<int> nums)
{
    std::sort(nums.begin(), nums.end());

    const int n = static_cast<int>(nums.size());
    std::vector<std::vector<int>> answer;
    std::vector<int> path;
    std::vector<char> used(n, false);

    struct State {
        std::vector<int> path;
        std::vector<char> used;
    };

    std::vector<State> stack{{path, used}};
    while (!stack.empty()) {
        State state = std::move(stack.back());
        stack.pop_back();

        if (static_cast<int>(state.path.size()) == n) {
            answer.push_back(std::move(state.path));
            continue;
        }

        for (int i = n - 1; i >= 0; --i) {
            if (state.used[i]) {
                continue;
            }
            if (i > 0 && nums[i] == nums[i - 1] && !state.used[i - 1]) {
                continue;
            }

            State next = state;
            next.used[i] = true;
            next.path.push_back(nums[i]);
            stack.push_back(std::move(next));
        }
    }

    return answer;
}
\end{lstlisting}

这个版本为了让去重条件清楚，直接把状态复制进栈，适合教材解释；高性能版本会用一份路径和标记做原地回退。去重条件最容易写反。排序后，相同数字挨在一起，\code{!state.used[i - 1]} 表示前一个相同数字在当前层还没有被拿来作为代表，此时选择后一个会制造重复分支，所以跳过。时间复杂度仍然与唯一排列数量成正比，最坏不超过 \complexity{O(n \cdot n!)}。

子集题的决策树更简单。每个元素只有两种选择：选或者不选。暴力枚举可以用二进制掩码，从 \code{0} 到 \code{2^n - 1}，第 \code{i} 位表示是否选择第 \code{i} 个元素。回溯写法则是沿着数组位置向后决定。两种方法本质相同，都是遍历大小为 \complexity{2^n} 的搜索空间。

\begin{lstlisting}[language=C++,caption={LeetCode 78 子集：迭代扩展已有答案}]
std::vector<std::vector<int>> subsets(const std::vector<int>& nums)
{
    std::vector<std::vector<int>> answer{{}};

    for (const int x : nums) {
        const int old_size = static_cast<int>(answer.size());
        for (int i = 0; i < old_size; ++i) {
            std::vector<int> next = answer[i];
            next.push_back(x);
            answer.push_back(std::move(next));
        }
    }

    return answer;
}
\end{lstlisting}

这段代码的不变量是：处理完前 \code{i} 个元素后，\code{answer} 包含这些元素的全部子集。加入新元素 \code{x} 时，旧子集保持“不选 \code{x}”的版本；每个旧子集复制一份并加上 \code{x}，得到“选 \code{x}”的版本。这个算法没有剪枝，因为所有子集都是合法答案。时间复杂度 \complexity{O(n2^n)}，空间复杂度同样由输出规模主导。

组合总和比子集多了一个目标约束。题目给定候选数组和目标值，要求找出所有和为目标的组合，同一个数字可以重复使用。暴力思路是不断选数字，直到长度达到某个上限再检查和；问题是上限不好定，而且大量分支会在总和已经超过目标后继续扩展。优化来自两个事实：候选数字都是正数；排序后，如果当前候选已经大于剩余目标，后面的更大候选也不可能使用。

\begin{lstlisting}[language=C++,caption={LeetCode 39 组合总和：显式状态保存起点和剩余目标}]
std::vector<std::vector<int>> combination_sum(std::vector<int> candidates,
                                              int target)
{
    std::sort(candidates.begin(), candidates.end());

    struct State {
        int start = 0;
        int remaining = 0;
        std::vector<int> path;
    };

    std::vector<std::vector<int>> answer;
    std::vector<State> stack{{0, target, {}}};

    while (!stack.empty()) {
        State state = std::move(stack.back());
        stack.pop_back();

        if (state.remaining == 0) {
            answer.push_back(std::move(state.path));
            continue;
        }

        for (int i = static_cast<int>(candidates.size()) - 1;
             i >= state.start;
             --i) {
            if (candidates[i] > state.remaining) {
                continue;
            }
            State next = state;
            next.path.push_back(candidates[i]);
            next.remaining -= candidates[i];
            next.start = i;
            stack.push_back(std::move(next));
        }
    }

    return answer;
}
\end{lstlisting}

\code{start} 的语义很重要：组合不关心顺序，\code{[2,3,2]} 和 \code{[2,2,3]} 是同一个答案，因此后续只能从当前下标及其右侧继续选，避免同一组合以不同顺序重复出现。因为允许重复使用当前数字，下一层的 \code{start} 仍然是 \code{i}；如果题目变成每个数字只能用一次，就应设为 \code{i + 1}。这就是组合类题最核心的变体开关。复杂度取决于目标值和候选分布，不能简单写成多项式；面试里应说明它是指数级搜索，但剪枝能显著减少无效分支。

括号生成的关键不是枚举所有长度为 \code{2n} 的括号串再检查。那样有 \complexity{2^{2n}} 个候选，大多数不合法。合法前缀必须满足两个计数约束：左括号数量不能超过 \code{n}；任何前缀中右括号数量不能超过左括号数量。把这两个约束提前放进生成过程，就不会产生明显非法的前缀。

\begin{lstlisting}[language=C++,caption={LeetCode 22 括号生成：前缀合法性剪枝}]
std::vector<std::string> generate_parenthesis(int n)
{
    struct State {
        std::string text;
        int left = 0;
        int right = 0;
    };

    std::vector<std::string> answer;
    std::vector<State> stack{{"", 0, 0}};

    while (!stack.empty()) {
        State state = std::move(stack.back());
        stack.pop_back();

        if (static_cast<int>(state.text.size()) == 2 * n) {
            answer.push_back(std::move(state.text));
            continue;
        }

        if (state.right < state.left) {
            State next = state;
            next.text.push_back(')');
            ++next.right;
            stack.push_back(std::move(next));
        }
        if (state.left < n) {
            State next = state;
            next.text.push_back('(');
            ++next.left;
            stack.push_back(std::move(next));
        }
    }

    return answer;
}
\end{lstlisting}

这题的正确性可以用前缀不变量解释：任何时刻 \code{left <= n} 且 \code{right <= left}。如果 \code{right > left}，说明某个右括号没有左括号匹配，后面再加字符也无法修复这个非法前缀；如果 \code{left > n}，也已经超过题目允许。输出数量是第 \code{n} 个卡特兰数，算法复杂度至少与答案数量成正比。易错点是终止条件：必须长度达到 \code{2n} 才能加入答案，只看 \code{left == right} 会把短前缀误加进去。

单词搜索把回溯放进网格图。节点是格子，边是上下左右相邻，路径不能重复使用同一个格子。暴力想法是从每个格子出发尝试所有方向，最后检查路径字符串；优化点是边走边比较目标单词的当前位置，一旦字符不匹配立即停止。访问标记也必须属于当前路径，而不是全局永久标记，因为同一个格子可以出现在不同起点或不同路径中。

\begin{lstlisting}[language=C++,caption={LeetCode 79 单词搜索：显式栈保存路径访问状态}]
bool exist(std::vector<std::vector<char>>& board, const std::string& word)
{
    if (word.empty()) {
        return true;
    }
    if (board.empty() || board[0].empty()) {
        return false;
    }

    constexpr std::array<std::pair<int, int>, 4> DIRECTIONS{
        std::pair{1, 0}, std::pair{-1, 0},
        std::pair{0, 1}, std::pair{0, -1}
    };

    const int rows = static_cast<int>(board.size());
    const int cols = static_cast<int>(board[0].size());

    struct State {
        int row = 0;
        int col = 0;
        int index = 0;
        std::vector<std::vector<char>> visited;
    };

    for (int row = 0; row < rows; ++row) {
        for (int col = 0; col < cols; ++col) {
            if (board[row][col] != word[0]) {
                continue;
            }

            std::vector<std::vector<char>> visited(rows, std::vector<char>(cols, false));
            visited[row][col] = true;
            std::vector<State> stack{{row, col, 0, std::move(visited)}};

            while (!stack.empty()) {
                State state = std::move(stack.back());
                stack.pop_back();
                if (state.index + 1 == static_cast<int>(word.size())) {
                    return true;
                }

                for (const auto [dr, dc] : DIRECTIONS) {
                    const int next_row = state.row + dr;
                    const int next_col = state.col + dc;
                    const int next_index = state.index + 1;
                    const bool out_of_bounds =
                        next_row < 0 || next_row >= rows ||
                        next_col < 0 || next_col >= cols;
                    if (out_of_bounds || state.visited[next_row][next_col] ||
                        board[next_row][next_col] != word[next_index]) {
                        continue;
                    }

                    State next = state;
                    next.row = next_row;
                    next.col = next_col;
                    next.index = next_index;
                    next.visited[next_row][next_col] = true;
                    stack.push_back(std::move(next));
                }
            }
        }
    }

    return false;
}
\end{lstlisting}

这个教材版把 \code{visited} 放进状态里，语义最清楚，但会复制较多布尔矩阵；原地回退版性能更好。本题最坏复杂度接近 \complexity{O(mn \cdot 3^L)}：第一步最多从每个格子出发，之后每一步通常不能回到上一个格子，所以分支近似 3，\code{L} 是单词长度。常见错误是把访问标记设成全局，导致某个格子被一条失败路径占用后，其他合法路径无法使用它。

N 皇后是回溯剪枝的代表。暴力方法在 \complexity{n \times n} 棋盘上选 \code{n} 个格子，再检查行、列、对角线是否冲突，搜索空间巨大。优化方式是按行放皇后：第 \code{row} 层只决定第 \code{row} 行皇后放在哪一列。这样行冲突天然消失，只需要维护列、主对角线和副对角线是否被占用。

\begin{lstlisting}[language=C++,caption={LeetCode 51 N 皇后：按行放置并维护三类占用}]
std::vector<std::vector<std::string>> solve_n_queens(int n)
{
    struct State {
        int row = 0;
        std::vector<int> queens;
        std::vector<char> cols;
        std::vector<char> diag1;
        std::vector<char> diag2;
    };

    std::vector<std::vector<std::string>> answer;
    std::vector<State> stack{{0,
                              {},
                              std::vector<char>(n, false),
                              std::vector<char>(2 * n - 1, false),
                              std::vector<char>(2 * n - 1, false)}};

    while (!stack.empty()) {
        State state = std::move(stack.back());
        stack.pop_back();

        if (state.row == n) {
            std::vector<std::string> board(n, std::string(n, '.'));
            for (int row = 0; row < n; ++row) {
                board[row][state.queens[row]] = 'Q';
            }
            answer.push_back(std::move(board));
            continue;
        }

        for (int col = n - 1; col >= 0; --col) {
            const int main_diag = state.row - col + n - 1;
            const int anti_diag = state.row + col;
            if (state.cols[col] || state.diag1[main_diag] ||
                state.diag2[anti_diag]) {
                continue;
            }

            State next = state;
            next.queens.push_back(col);
            next.cols[col] = true;
            next.diag1[main_diag] = true;
            next.diag2[anti_diag] = true;
            ++next.row;
            stack.push_back(std::move(next));
        }
    }

    return answer;
}
\end{lstlisting}

主对角线上的格子满足 \code{row - col} 相同；为了把负数下标变成数组下标，代码加上 \code{n - 1}。副对角线满足 \code{row + col} 相同，范围天然是 \code{0} 到 \code{2n - 2}。这题的面试重点不是把棋盘画出来，而是把冲突检测从扫描整张棋盘优化成三个数组的常数查询。复杂度仍然是指数级，但剪枝把不可行分支尽早挡在下一层之外。

回溯去重可以分成两类：路径内去重和同层去重。路径内去重解决“同一个元素不能被重复使用”，典型工具是 \code{used} 数组；同层去重解决“相同值作为同一层选择会生成重复答案”，典型工具是排序后跳过相邻重复值。全排列 II 同时需要这两类去重，组合总和 II 主要需要同层去重。只要把这两个概念混在一起，条件就很容易写反。

组合总和 II 中，每个候选数字只能使用一次，输入可能有重复值，答案组合不能重复。与组合总和 I 相比，下一层起点必须是 \code{i + 1}，因为当前下标不能再次使用；与全排列 II 相比，我们不需要 \code{used} 数组，因为 \code{start} 已经保证每个下标最多被使用一次。同层去重条件是：在同一个循环层里，如果当前值和前一个值相同，并且前一个值已经作为本层候选被跳过或处理过，那么当前值也不能再作为新的分支开头。

\begin{lstlisting}[language=C++,caption={LeetCode 40 组合总和 II：每个下标最多使用一次，同层跳过重复值}]
std::vector<std::vector<int>> combination_sum2(std::vector<int> candidates,
                                               int target)
{
    std::sort(candidates.begin(), candidates.end());

    struct State {
        int start = 0;
        int remaining = 0;
        std::vector<int> path;
    };

    std::vector<std::vector<int>> answer;
    std::vector<State> stack{{0, target, {}}};

    while (!stack.empty()) {
        State state = std::move(stack.back());
        stack.pop_back();

        if (state.remaining == 0) {
            answer.push_back(std::move(state.path));
            continue;
        }

        for (int i = static_cast<int>(candidates.size()) - 1;
             i >= state.start;
             --i) {
            if (i + 1 < static_cast<int>(candidates.size()) &&
                candidates[i] == candidates[i + 1] &&
                i + 1 >= state.start) {
                continue;
            }
            if (candidates[i] > state.remaining) {
                continue;
            }

            State next = state;
            next.path.push_back(candidates[i]);
            next.remaining -= candidates[i];
            next.start = i + 1;
            stack.push_back(std::move(next));
        }
    }

    return answer;
}
\end{lstlisting}

这段显式栈代码为了保持输出方向，用倒序枚举候选；倒序时同层去重要和相邻的后一个元素比较。递归版本通常正序枚举，条件写成 \code{i > start && candidates[i] == candidates[i - 1]}。两个条件本质相同：同一个递归层里，相同值只能由一个下标代表。面试中如果不确定倒序条件，先写递归式正序伪代码说明去重，再实现自己更熟悉的版本。最重要的是讲清楚：同层重复会制造相同组合，路径重复则是同一个下标重复使用，这两个问题不能用同一个条件糊过去。

字符串切分题也属于回溯。分割回文串要求把字符串切成若干段，每一段都是回文。暴力做法是枚举所有切分方式，再逐段检查是否回文；如果每次检查都扫描子串，会产生大量重复判断。优化方法是先用 DP 预处理 \code{is_palindrome[left][right]}，表示闭区间子串是否回文。之后回溯只需要常数时间判断某一段是否合法。

\begin{lstlisting}[language=C++,caption={LeetCode 131 分割回文串：预处理回文表后搜索切分点}]
std::vector<std::vector<std::string>> partition_palindrome(const std::string& s)
{
    const int n = static_cast<int>(s.size());
    std::vector<std::vector<char>> is_palindrome(n, std::vector<char>(n, false));
    for (int length = 1; length <= n; ++length) {
        for (int left = 0; left + length - 1 < n; ++left) {
            const int right = left + length - 1;
            if (s[left] != s[right]) {
                continue;
            }
            is_palindrome[left][right] =
                length <= 2 || is_palindrome[left + 1][right - 1];
        }
    }

    struct State {
        int start = 0;
        std::vector<std::string> parts;
    };

    std::vector<std::vector<std::string>> answer;
    std::vector<State> stack{{0, {}}};
    while (!stack.empty()) {
        State state = std::move(stack.back());
        stack.pop_back();
        if (state.start == n) {
            answer.push_back(std::move(state.parts));
            continue;
        }

        for (int end = n - 1; end >= state.start; --end) {
            if (!is_palindrome[state.start][end]) {
                continue;
            }
            State next = state;
            next.parts.push_back(s.substr(state.start, end - state.start + 1));
            next.start = end + 1;
            stack.push_back(std::move(next));
        }
    }
    return answer;
}
\end{lstlisting}

这题的决策树层数不是固定的 \code{n}，而是取决于切了多少段。每个状态的 \code{start} 表示下一个待切分位置，选择是从 \code{start} 开始切到哪个 \code{end}。预处理回文表的转移也很经典：两端字符相等，且内部子串是回文，整个子串才是回文；长度为 1 或 2 时内部为空或单字符，可以直接成立。搜索复杂度仍然可能指数级，因为合法切分数量本身可能很多，但预处理避免了在搜索树中反复扫描同一段字符串。

复原 IP 地址是切分题的约束版。字符串要切成四段，每段长度 1 到 3，数值 0 到 255，且不能有前导零。暴力枚举三个切分点就能做，因为段数固定；也可以用回溯统一表达。剪枝点很明确：剩余字符太少或太多时停止；当前段长度超过 3 停止；以 0 开头的多位段非法；数值超过 255 非法。

\begin{lstlisting}[language=C++,caption={LeetCode 93 复原 IP 地址：段数和剩余长度剪枝}]
std::vector<std::string> restore_ip_addresses(const std::string& s)
{
    constexpr int SEGMENT_COUNT = 4;
    constexpr int MAX_SEGMENT_VALUE = 255;
    std::vector<std::string> answer;

    struct State {
        int index = 0;
        std::vector<std::string> segments;
    };

    std::vector<State> stack{{0, {}}};
    while (!stack.empty()) {
        State state = std::move(stack.back());
        stack.pop_back();

        const int used_segments = static_cast<int>(state.segments.size());
        const int remaining_segments = SEGMENT_COUNT - used_segments;
        const int remaining_chars = static_cast<int>(s.size()) - state.index;
        if (remaining_chars < remaining_segments ||
            remaining_chars > remaining_segments * 3) {
            continue;
        }

        if (used_segments == SEGMENT_COUNT) {
            if (state.index == static_cast<int>(s.size())) {
                answer.push_back(state.segments[0] + "." + state.segments[1] +
                                 "." + state.segments[2] + "." +
                                 state.segments[3]);
            }
            continue;
        }

        int value = 0;
        for (int length = 1; length <= 3; ++length) {
            if (state.index + length > static_cast<int>(s.size())) {
                break;
            }
            if (length > 1 && s[state.index] == '0') {
                break;
            }
            value = value * 10 + (s[state.index + length - 1] - '0');
            if (value > MAX_SEGMENT_VALUE) {
                break;
            }

            State next = state;
            next.segments.push_back(s.substr(state.index, length));
            next.index += length;
            stack.push_back(std::move(next));
        }
    }

    return answer;
}
\end{lstlisting}

这题适合训练“剩余量剪枝”。如果剩余字符数少于剩余段数，每段至少一个字符都不够；如果剩余字符数大于剩余段数乘 3，每段最多三个字符也放不下。这种剪枝不依赖数据运气，是严格必要条件，因此不会漏解。很多组合题都可以做类似估算：剩余元素数量是否足够，剩余目标是否可能达到，剩余步数是否还能完成。

回溯题也可以用位掩码表示集合。子集枚举、访问状态、N 皇后列占用都可以压成整数位。位掩码的优势是状态复制便宜、哈希和比较方便；缺点是可读性下降，并且 \code{n} 不能太大。对于 \code{n <= 20} 的集合问题，位掩码常常是从回溯过渡到状态压缩 DP 的桥梁。

\begin{lstlisting}[language=C++,caption={位掩码枚举子集：每一位表示选或不选}]
std::vector<std::vector<int>> subsets_by_mask(const std::vector<int>& nums)
{
    const int n = static_cast<int>(nums.size());
    const int state_count = 1 << n;
    std::vector<std::vector<int>> answer;
    answer.reserve(state_count);

    for (int mask = 0; mask < state_count; ++mask) {
        std::vector<int> subset;
        for (int i = 0; i < n; ++i) {
            if ((mask & (1 << i)) != 0) {
                subset.push_back(nums[i]);
            }
        }
        answer.push_back(std::move(subset));
    }

    return answer;
}
\end{lstlisting}

位掩码写法的复杂度是 \complexity{O(n2^n)}，因为每个集合要检查 n 位。它没有递归栈，也不需要回退，适合子集这种所有状态都合法的题；如果有复杂剪枝，例如组合总和或单词搜索，回溯更自然。选择表达方式时要看题目结构，不要为了使用位运算牺牲清晰性。

剪枝必须证明不会漏解。常见剪枝有四类。第一，合法性剪枝：括号前缀右括号不能超过左括号。第二，容量剪枝：组合总和中当前和超过目标立即停止。第三，剩余量剪枝：IP 地址剩余字符数量必须落在可行范围。第四，排序剪枝：候选已排序时，当前值超过剩余目标，后续更大值也不可能使用。没有证明的剪枝很危险，比如“看起来这个分支不太可能有答案”不能作为正确算法的一部分。

\begin{labbox}
回溯实验：第一，把组合总和 II 的同层去重条件删除，观察重复答案如何出现。第二，把复原 IP 地址的剩余量剪枝删除，统计访问状态数量变化。第三，把分割回文串的回文预处理去掉，改成每次扫描判断，比较长字符串上的耗时。第四，用位掩码和迭代扩展两种方式生成子集，比较输出顺序和代码复杂度。
\end{labbox}

\begin{principlebox}{回溯题复盘模板}
回溯题不要先背代码。先画出决策树的一层，写清本层选择、路径状态、可选集合、终止条件和剪枝条件。组合题重点看 \code{start} 是否允许重复使用当前元素；排列题重点看 \code{used} 和同层去重；网格题重点看访问标记是路径级还是全局级；括号和棋盘题重点把非法前缀提前剪掉。最后再问一句：答案数量本身有多大？如果答案就是指数级，复杂度不可能被写成线性或多项式。
\end{principlebox}

回溯、枚举和剪枝深度题卡 按题型拆成章内大节，但每个大节都先从 LeetCode 案例切入，再进入分流、案例矩阵、案例推导和实验复盘。这样读者看到的是从具体题推到规律，而不是先读抽象方法再猜它能用在哪。

\topic{回溯、枚举与剪枝}
这一节先看 LeetCode 17 电话号码的字母组合、LeetCode 22 括号生成、LeetCode 39 组合总和 这几道 LeetCode 题，再整理同一题型的分流和推导。阅读顺序不是先背原理，而是先看案例的暴力瓶颈，再把重复出现的观察抽象成方法。

\subsection{原理和适用条件}

以 LeetCode 17 电话号码的字母组合、LeetCode 22 括号生成、LeetCode 39 组合总和 为主线：LeetCode 17 电话号码的字母组合：模型是“每个数字对应若干字母，答案是所有按位选择形成的字符串。”，优化抓手是“暴力就是完整笛卡尔积，优化重点是路径长度和下标同步推进，避免在每层重新拼接大量中间字符串。”；LeetCode 22 括号生成：模型是“合法括号串的前缀中右括号数量不能超过左括号数量。”，优化抓手是“状态保存已用左括号和右括号数量，只有满足约束的选择才继续递归。”；LeetCode 39 组合总和：模型是“候选数字可重复使用，答案不关心排列顺序。”，优化抓手是“排序后按起点枚举，递归时仍传当前下标表示可以继续使用同一数字。”。这些案例共同推出的规律是：先写搜索树：路径是什么、选择列表是什么、什么时候结束。然后用排序、剩余容量、同层去重或合法性预处理剪掉不可能分支。 方法名只是复盘后的标签，真正要掌握的是从题面约束、暴力失败、状态设计到边界反例的推导链。

\begin{principlebox}{图示：从 LeetCode 案例到 回溯、枚举与剪枝推导链}
\begin{tabular}{@{}p{0.18\linewidth}p{0.74\linewidth}@{}}
暴力基线 & 枚举候选、完整模拟或逐个检查，先保证覆盖全部可能答案。\\
瓶颈定位 & 慢点在于大量分支在很早就已经不可能成功，仍然被继续展开。重复元素还会产生等价排列或组合。\\
结构选择 & 剪枝来自容量、剩余长度、排序后去重、合法性预检查和提前终止。组合题关注起点推进，排列题关注元素是否使用，分割题关注当前片段是否合法。\\
不变量 & 路径数组保存已经做出的选择；起点下标保证组合不回头；同层去重保证同一个位置层级不会选择相同值；已使用数组保证排列不重复使用同一元素。\\
代码落地 & 回溯代码虽然天然递归，但工程中要意识到深度。题目规模较小时递归可读性好；若要遵守严格工程栈限制，可以改成显式栈状态机。\\
\end{tabular}
\end{principlebox}

\subsection{LeetCode 案例分流}

\begin{principlebox}{从 LeetCode 案例分流：回溯、枚举与剪枝}
\textbf{子集和组合。}
代表案例：LeetCode 78 子集、LeetCode 90 子集 II、LeetCode 77 组合、LeetCode 216 组合总和 III。先看这些题的题面条件：答案不关心顺序，元素按起点向后选择。由此推出的处理方式是：递增起点防止重复排列，同层去重防止重复值重复生成。风险点：同层去重和路径去重不是一回事。

\textbf{排列。}
代表案例：LeetCode 46 全排列、LeetCode 47 全排列 II。先看这些题的题面条件：答案关心顺序，每个位置选择一个未使用元素。由此推出的处理方式是：used 数组记录路径中已使用对象。风险点：重复元素排列要排序并跳过同层等价选择。

\textbf{分割和构造。}
代表案例：LeetCode 131 分割回文串、LeetCode 93 复原 IP 地址、LeetCode 140 单词拆分 II。先看这些题的题面条件：每一步选择一个前缀、片段或局部结构。由此推出的处理方式是：检查片段合法性后进入下一位置。风险点：输出规模可能指数级，复杂度要包含输出。

\textbf{棋盘和约束搜索。}
代表案例：LeetCode 51 N 皇后、LeetCode 79 单词搜索、LeetCode 212 单词搜索 II。先看这些题的题面条件：选择会影响行列、对角线或局部访问状态。由此推出的处理方式是：用集合或位掩码保存约束，进入退出时成对修改。风险点：剪枝只能删除不可能成功的分支。

\end{principlebox}

\subsection{案例矩阵}

本节案例覆盖 LeetCode 17 电话号码的字母组合、LeetCode 22 括号生成、LeetCode 39 组合总和、LeetCode 40 组合总和 II、LeetCode 46 全排列、LeetCode 47 全排列 II、LeetCode 51 N 皇后、LeetCode 78 子集、LeetCode 79 单词搜索、LeetCode 90 子集 II、LeetCode 93 复原 IP 地址、LeetCode 131 分割回文串、LeetCode 140 单词拆分 II、LeetCode 212 单词搜索 II、LeetCode 216 组合总和 III、LeetCode 679 24 点游戏、LeetCode 698 划分为 k 个相等的子集。阅读时不要按题号背答案，而要先判断题面落在哪个分流场景：查询目标是什么，输入约束给了什么单调性或等价关系，暴力慢在哪里，优化结构保存了什么状态。

\begin{principlebox}{案例矩阵}
\textbf{LeetCode 17 电话号码的字母组合。}
每个数字对应若干字母，答案是所有按位选择形成的字符串。单点风险：空输入、包含 7 或 9、输入中不应出现 0 和 1、输出规模指数增长。

\textbf{LeetCode 22 括号生成。}
合法括号串的前缀中右括号数量不能超过左括号数量。单点风险：n 为零或一、右括号提前超过左括号、输出规模是卡特兰数。

\textbf{LeetCode 39 组合总和。}
候选数字可重复使用，答案不关心排列顺序。单点风险：目标小于最小数、刚好等于、候选有重复时题意如何处理、输出规模。

\textbf{LeetCode 40 组合总和 II。}
每个候选只能使用一次，输入中可能存在重复值。单点风险：全重复、目标为零、同层去重写成跨层去重、答案顺序不要求。

\textbf{LeetCode 46 全排列。}
排列关心元素顺序，每个位置从所有未使用元素中选择一个。单点风险：空数组、单元素、输出规模为 n 阶乘、元素值唯一假设。

\textbf{LeetCode 47 全排列 II。}
重复元素会让不同下标选择生成相同排列。单点风险：全重复、部分重复、同层去重条件写反、输出顺序不要求。

\textbf{LeetCode 51 N 皇后。}
每一行放一个皇后，列和两条对角线不能冲突。单点风险：n 为一、n 为二或三无解、对角线编号、回退时状态恢复。

\textbf{LeetCode 78 子集。}
每个元素都有选或不选两种选择，答案是所有路径节点。单点风险：空数组、单元素、输出数量为二的 n 次方、顺序不要求。

\textbf{LeetCode 79 单词搜索。}
网格路径不能重复使用同一格，状态包含位置、匹配下标和访问标记。单点风险：单字符单词、路径需要回退、重复字母、单元格不能复用。

\textbf{LeetCode 90 子集 II。}
重复元素会让不同下标路径生成相同集合。单点风险：全重复、空集、重复值相邻、去重条件写成同层而不是全局。

\textbf{LeetCode 93 复原 IP 地址。}
字符串要切成四段，每段必须是 0 到 255 且无非法前导零。单点风险：前导零、数值超过 255、剩余字符太多或太少、刚好四段。

\textbf{LeetCode 131 分割回文串。}
每次选择一个从当前位置开始的回文前缀。单点风险：空串、单字符、全相同字符、重复分割路径。

\textbf{LeetCode 140 单词拆分 II。}
需要输出所有由字典单词拼出的句子，输出规模本身可能指数级。单点风险：无解后缀、重复单词、路径拼接成本、答案数量主导复杂度。

\textbf{LeetCode 212 单词搜索 II。}
多个单词在同一棋盘上搜索，公共前缀可以共享。单点风险：重复单词、同一格不能复用、前缀是完整单词、剪枝删除空子树。

\textbf{LeetCode 216 组合总和 III。}
从一到九中选 k 个数和为 n，每个数最多用一次。单点风险：n 太小、n 太大、k 为零、路径长度已经超过 k。

\textbf{LeetCode 679 24 点游戏。}
四个数字必须全部使用，括号表达式可以看成每次选两个数合并成一个新数。单点风险：减法和除法顺序、除零、浮点误差、重复中间状态。

\textbf{LeetCode 698 划分为 k 个相等的子集。}
每个元素要分配到 k 个桶中，所有桶目标和相同。单点风险：存在元素大于目标、重复元素、空桶对称、k 为一。

\end{principlebox}

\subsection{共性落点}

\begin{principlebox}{本组共性落点}
\textbf{慢版本。} 暴力回溯会枚举整个搜索树。它不是错误，而是必须先写清楚选择列表、路径状态和终止条件，再讨论剪枝。
\textbf{瓶颈。} 慢点在于大量分支在很早就已经不可能成功，仍然被继续展开。重复元素还会产生等价排列或组合。
\textbf{状态字段。} 路径、起点或使用标记、剩余目标、答案集合、剪枝条件；进入递归和退出递归必须成对恢复。
\textbf{不变量。} 路径数组保存已经做出的选择；起点下标保证组合不回头；同层去重保证同一个位置层级不会选择相同值；已使用数组保证排列不重复使用同一元素。
\textbf{实现检查。} 剪枝条件写在扩展前；同层去重不要误写成全局去重；保存答案时复制当前路径。
\textbf{C++ 落地。} 回溯代码虽然天然递归，但工程中要意识到深度。题目规模较小时递归可读性好；若要遵守严格工程栈限制，可以改成显式栈状态机。
\textbf{证明抓手。} 证明从搜索树覆盖开始：每个合法答案有且只有一条路径，剪枝只删掉不可能成功的分支。
\end{principlebox}

\subsection{案例推导}

\noindent\textbf{LeetCode 17 电话号码的字母组合。}

\textbf{题目说明。} LeetCode 17 电话号码的字母组合 的题面入口要先还原成具体任务：每个数字对应若干字母，答案是所有按位选择形成的字符串。

\textbf{样例入口。} 拿 digits=23 手算，2 对应 a/b/c，3 对应 d/e/f，第一层选 a 后第二层依次选 d/e/f，得到 ad、ae、af。

\textbf{暴力基线。} 慢版本就是完整搜索树：每层列出选择列表，路径到达终止条件时检查答案。它先保证覆盖所有可能，再把明显无效或重复的分支剪掉。放到本题就是：先按“每个数字对应若干字母，答案是所有按位选择形成的字符串”完整试慢版本；样例里已经能看到“规律是层数对应输入数字位置，路径保存已经选择的字母；空输入应返回空结果而不是包含空串”，所以优化前必须先能说清慢版本枚举了哪些候选。

\textbf{暴力过程。} 拿 digits=23 手算，2 对应 a/b/c，3 对应 d/e/f，第一层选 a 后第二层依次选 d/e/f，得到 ad、ae、af。

\textbf{重复工作。} 题面模型：每个数字对应若干字母，答案是所有按位选择形成的字符串。暴力版的慢点要落到本题：慢版本会围绕“每个数字对应若干字母，答案是所有按位选择形成的字符串”展开完整搜索树。样例规律是“规律是层数对应输入数字位置，路径保存已经选择的字母；空输入应返回空结果而不是包含空串”，说明一层递归必须对应一个明确决策，剪枝只能删除整棵不可能成功或重复的子树。后面的优化只消掉这类重复工作，不能改变答案集合。

\textbf{优化条件。} 本题的关键观察是：暴力就是完整笛卡尔积，优化重点是路径长度和下标同步推进，避免在每层重新拼接大量中间字符串。这句话必须能翻译成可维护状态；如果题目条件改变后观察不再成立，就不能继续套原来的写法。

\textbf{相邻状态。} 规律是层数对应输入数字位置，路径保存已经选择的字母；空输入应返回空结果而不是包含空串。把这个特例继续拆开看：一层递归对应一个明确决策，路径保存已经做出的选择，选择列表保存仍可能扩展的分支。剪枝前必须能说明被剪掉的整棵子树没有合法答案，而不是只因为它看起来不优。

\textbf{状态复用。} 把完整搜索树加上合法性检查、剩余资源剪枝和同层去重；路径状态进入和退出要成对恢复。本题落点是：规律是层数对应输入数字位置，路径保存已经选择的字母；空输入应返回空结果而不是包含空串。手算时只改被新输入影响的那一小部分，而不是回头重算完整候选。

\textbf{指针和字段。} 状态字段要能说清：路径、起点或使用标记、剩余目标、答案集合、剪枝条件；进入递归和退出递归必须成对恢复。手算时按这个协议跟踪：拿 digits=23 手算，2 对应 a/b/c，3 对应 d/e/f，第一层选 a 后第二层依次选 d/e/f，得到 ad、ae、af。然后按协议复查：手算时给每层标出选择列表和路径。剪枝发生前要指出被剪分支为什么已经不可能得到合法答案。

\textbf{代码落点。} 剪枝条件写在扩展前；同层去重不要误写成全局去重；保存答案时复制当前路径。关键代码行必须能逐句翻译成“旧状态怎样变成新状态”，否则就只是模板。

\textbf{正确性。} 证明从搜索树覆盖开始：每个合法答案有且只有一条路径，剪枝只删掉不可能成功的分支。

\textbf{边界实验。} 先用空输入、包含 7 或 9、输入中不应出现 0 和 1、输出规模指数增长做反例检查。实验时覆盖空输入、包含 7 或 9、输入中不应出现 0 和 1、输出规模指数增长，统计搜索节点数量和剪枝次数。重复元素题要打印同一层跳过哪些值，确认没有误删合法路径。

\textbf{复杂度变化。} 暴力搜索树的规模由输出或选择空间决定；剪枝不改变最坏指数级本质，但能删除不可能成功或重复的分支，复杂度要说明输出规模。

\textbf{迁移追问。} 如果题目改成“若字符映射来自键盘以外的字典，只需要替换候选表，搜索骨架不变”，先判断原状态是否仍然覆盖新答案集合，再决定是微调边界、改 value 语义，还是换算法。

\noindent\textbf{LeetCode 22 括号生成。}

\textbf{题目说明。} LeetCode 22 括号生成 的题面入口要先还原成具体任务：合法括号串的前缀中右括号数量不能超过左括号数量。

\textbf{样例入口。} 拿 n=2 手算，第一步只能放左括号，之后可以放左括号得到 (())，也可以在已有左括号数量大于右括号时放右括号得到 ()()。

\textbf{暴力基线。} 慢版本就是完整搜索树：每层列出选择列表，路径到达终止条件时检查答案。它先保证覆盖所有可能，再把明显无效或重复的分支剪掉。放到本题就是：先按“合法括号串的前缀中右括号数量不能超过左括号数量”完整试慢版本；样例里已经能看到“规律是状态包含已放左括号数和右括号数，右括号不能超过左括号，左括号不能超过 n”，所以优化前必须先能说清慢版本枚举了哪些候选。

\textbf{暴力过程。} 拿 n=2 手算，第一步只能放左括号，之后可以放左括号得到 (())，也可以在已有左括号数量大于右括号时放右括号得到 ()()。

\textbf{重复工作。} 题面模型：合法括号串的前缀中右括号数量不能超过左括号数量。暴力版的慢点要落到本题：慢版本会围绕“合法括号串的前缀中右括号数量不能超过左括号数量”展开完整搜索树。样例规律是“规律是状态包含已放左括号数和右括号数，右括号不能超过左括号，左括号不能超过 n”，说明一层递归必须对应一个明确决策，剪枝只能删除整棵不可能成功或重复的子树。后面的优化只消掉这类重复工作，不能改变答案集合。

\textbf{优化条件。} 本题的关键观察是：状态保存已用左括号和右括号数量，只有满足约束的选择才继续递归。这句话必须能翻译成可维护状态；如果题目条件改变后观察不再成立，就不能继续套原来的写法。

\textbf{相邻状态。} 规律是状态包含已放左括号数和右括号数，右括号不能超过左括号，左括号不能超过 n。把这个特例继续拆开看：一层递归对应一个明确决策，路径保存已经做出的选择，选择列表保存仍可能扩展的分支。剪枝前必须能说明被剪掉的整棵子树没有合法答案，而不是只因为它看起来不优。

\textbf{状态复用。} 把完整搜索树加上合法性检查、剩余资源剪枝和同层去重；路径状态进入和退出要成对恢复。本题落点是：规律是状态包含已放左括号数和右括号数，右括号不能超过左括号，左括号不能超过 n。手算时只改被新输入影响的那一小部分，而不是回头重算完整候选。

\textbf{指针和字段。} 状态字段要能说清：路径、起点或使用标记、剩余目标、答案集合、剪枝条件；进入递归和退出递归必须成对恢复。手算时按这个协议跟踪：拿 n=2 手算，第一步只能放左括号，之后可以放左括号得到 (())，也可以在已有左括号数量大于右括号时放右括号得到 ()()。然后按协议复查：手算时给每层标出选择列表和路径。剪枝发生前要指出被剪分支为什么已经不可能得到合法答案。

\textbf{代码落点。} 剪枝条件写在扩展前；同层去重不要误写成全局去重；保存答案时复制当前路径。关键代码行必须能逐句翻译成“旧状态怎样变成新状态”，否则就只是模板。

\textbf{正确性。} 证明从搜索树覆盖开始：每个合法答案有且只有一条路径，剪枝只删掉不可能成功的分支。

\textbf{边界实验。} 先用n 为零或一、右括号提前超过左括号、输出规模是卡特兰数做反例检查。实验时覆盖n 为零或一、右括号提前超过左括号、输出规模是卡特兰数，统计搜索节点数量和剪枝次数。重复元素题要打印同一层跳过哪些值，确认没有误删合法路径。

\textbf{复杂度变化。} 暴力搜索树的规模由输出或选择空间决定；剪枝不改变最坏指数级本质，但能删除不可能成功或重复的分支，复杂度要说明输出规模。

\textbf{迁移追问。} 如果题目改成“若有多种括号类型，状态还需要栈保存类型匹配关系”，先判断原状态是否仍然覆盖新答案集合，再决定是微调边界、改 value 语义，还是换算法。

\noindent\textbf{LeetCode 39 组合总和。}

\textbf{题目说明。} LeetCode 39 组合总和 的题面入口要先还原成具体任务：候选数字可重复使用，答案不关心排列顺序。

\textbf{样例入口。} 拿 candidates=[2, 3, 6, 7]、target=7 手算，选择 2 后仍可以继续选择 2，因为每个数可重复；路径 2, 2, 3 成功，7 单独也成功。

\textbf{暴力基线。} 慢版本就是完整搜索树：每层列出选择列表，路径到达终止条件时检查答案。它先保证覆盖所有可能，再把明显无效或重复的分支剪掉。放到本题就是：先按“候选数字可重复使用，答案不关心排列顺序”完整试慢版本；样例里已经能看到“规律是回溯起点不回退，但递归下一层仍传当前下标以允许重复使用”，所以优化前必须先能说清慢版本枚举了哪些候选。

\textbf{暴力过程。} 拿 candidates=[2, 3, 6, 7]、target=7 手算，选择 2 后仍可以继续选择 2，因为每个数可重复；路径 2, 2, 3 成功，7 单独也成功。

\textbf{重复工作。} 题面模型：候选数字可重复使用，答案不关心排列顺序。暴力版的慢点要落到本题：慢版本会围绕“候选数字可重复使用，答案不关心排列顺序”展开完整搜索树。样例规律是“规律是回溯起点不回退，但递归下一层仍传当前下标以允许重复使用”，说明一层递归必须对应一个明确决策，剪枝只能删除整棵不可能成功或重复的子树。后面的优化只消掉这类重复工作，不能改变答案集合。

\textbf{优化条件。} 本题的关键观察是：排序后按起点枚举，递归时仍传当前下标表示可以继续使用同一数字。这句话必须能翻译成可维护状态；如果题目条件改变后观察不再成立，就不能继续套原来的写法。

\textbf{相邻状态。} 规律是回溯起点不回退，但递归下一层仍传当前下标以允许重复使用。把这个特例继续拆开看：一层递归对应一个明确决策，路径保存已经做出的选择，选择列表保存仍可能扩展的分支。剪枝前必须能说明被剪掉的整棵子树没有合法答案，而不是只因为它看起来不优。

\textbf{状态复用。} 把完整搜索树加上合法性检查、剩余资源剪枝和同层去重；路径状态进入和退出要成对恢复。本题落点是：规律是回溯起点不回退，但递归下一层仍传当前下标以允许重复使用。手算时只改被新输入影响的那一小部分，而不是回头重算完整候选。

\textbf{指针和字段。} 状态字段要能说清：路径、起点或使用标记、剩余目标、答案集合、剪枝条件；进入递归和退出递归必须成对恢复。手算时按这个协议跟踪：拿 candidates=[2, 3, 6, 7]、target=7 手算，选择 2 后仍可以继续选择 2，因为每个数可重复；路径 2, 2, 3 成功，7 单独也成功。然后按协议复查：手算时给每层标出选择列表和路径。剪枝发生前要指出被剪分支为什么已经不可能得到合法答案。

\textbf{代码落点。} 剪枝条件写在扩展前；同层去重不要误写成全局去重；保存答案时复制当前路径。关键代码行必须能逐句翻译成“旧状态怎样变成新状态”，否则就只是模板。

\textbf{正确性。} 证明从搜索树覆盖开始：每个合法答案有且只有一条路径，剪枝只删掉不可能成功的分支。

\textbf{边界实验。} 先用目标小于最小数、刚好等于、候选有重复时题意如何处理、输出规模做反例检查。实验时覆盖目标小于最小数、刚好等于、候选有重复时题意如何处理、输出规模，统计搜索节点数量和剪枝次数。重复元素题要打印同一层跳过哪些值，确认没有误删合法路径。

\textbf{复杂度变化。} 暴力搜索树的规模由输出或选择空间决定；剪枝不改变最坏指数级本质，但能删除不可能成功或重复的分支，复杂度要说明输出规模。

\textbf{迁移追问。} 如果题目改成“组合总和 II 每个数只能用一次，起点推进和去重条件都不同”，先判断原状态是否仍然覆盖新答案集合，再决定是微调边界、改 value 语义，还是换算法。

\noindent\textbf{LeetCode 40 组合总和 II。}

\textbf{题目说明。} LeetCode 40 组合总和 II 的题面入口要先还原成具体任务：每个候选只能使用一次，输入中可能存在重复值。

\textbf{样例入口。} 拿 candidates=[10, 1, 2, 7, 6, 1, 5]、target=8 手算，排序后两个 1 在同一层只能选第一个作为起点，否则会生成重复组合；但选了第一个 1 后，下一层可以选第二个 1。

\textbf{暴力基线。} 慢版本就是完整搜索树：每层列出选择列表，路径到达终止条件时检查答案。它先保证覆盖所有可能，再把明显无效或重复的分支剪掉。放到本题就是：先按“每个候选只能使用一次，输入中可能存在重复值”完整试慢版本；样例里已经能看到“规律是每个数只能用一次，递归下一层传 i+1，并做同层去重”，所以优化前必须先能说清慢版本枚举了哪些候选。

\textbf{暴力过程。} 拿 candidates=[10, 1, 2, 7, 6, 1, 5]、target=8 手算，排序后两个 1 在同一层只能选第一个作为起点，否则会生成重复组合；但选了第一个 1 后，下一层可以选第二个 1。

\textbf{重复工作。} 题面模型：每个候选只能使用一次，输入中可能存在重复值。暴力版的慢点要落到本题：慢版本会围绕“每个候选只能使用一次，输入中可能存在重复值”展开完整搜索树。样例规律是“规律是每个数只能用一次，递归下一层传 i+1，并做同层去重”，说明一层递归必须对应一个明确决策，剪枝只能删除整棵不可能成功或重复的子树。后面的优化只消掉这类重复工作，不能改变答案集合。

\textbf{优化条件。} 本题的关键观察是：排序后递归下一层从 i 加一开始，同层遇到相同值要跳过。这句话必须能翻译成可维护状态；如果题目条件改变后观察不再成立，就不能继续套原来的写法。

\textbf{相邻状态。} 规律是每个数只能用一次，递归下一层传 i+1，并做同层去重。把这个特例继续拆开看：一层递归对应一个明确决策，路径保存已经做出的选择，选择列表保存仍可能扩展的分支。剪枝前必须能说明被剪掉的整棵子树没有合法答案，而不是只因为它看起来不优。

\textbf{状态复用。} 把完整搜索树加上合法性检查、剩余资源剪枝和同层去重；路径状态进入和退出要成对恢复。本题落点是：规律是每个数只能用一次，递归下一层传 i+1，并做同层去重。手算时只改被新输入影响的那一小部分，而不是回头重算完整候选。

\textbf{指针和字段。} 状态字段要能说清：路径、起点或使用标记、剩余目标、答案集合、剪枝条件；进入递归和退出递归必须成对恢复。手算时按这个协议跟踪：拿 candidates=[10, 1, 2, 7, 6, 1, 5]、target=8 手算，排序后两个 1 在同一层只能选第一个作为起点，否则会生成重复组合；但选了第一个 1 后，下一层可以选第二个 1。然后按协议复查：手算时给每层标出选择列表和路径。剪枝发生前要指出被剪分支为什么已经不可能得到合法答案。

\textbf{代码落点。} 剪枝条件写在扩展前；同层去重不要误写成全局去重；保存答案时复制当前路径。关键代码行必须能逐句翻译成“旧状态怎样变成新状态”，否则就只是模板。

\textbf{正确性。} 证明从搜索树覆盖开始：每个合法答案有且只有一条路径，剪枝只删掉不可能成功的分支。

\textbf{边界实验。} 先用全重复、目标为零、同层去重写成跨层去重、答案顺序不要求做反例检查。实验时覆盖全重复、目标为零、同层去重写成跨层去重、答案顺序不要求，统计搜索节点数量和剪枝次数。重复元素题要打印同一层跳过哪些值，确认没有误删合法路径。

\textbf{复杂度变化。} 暴力搜索树的规模由输出或选择空间决定；剪枝不改变最坏指数级本质，但能删除不可能成功或重复的分支，复杂度要说明输出规模。

\textbf{迁移追问。} 如果题目改成“子集 II 使用同样的同层去重，只是每个路径节点都可能成为答案”，先判断原状态是否仍然覆盖新答案集合，再决定是微调边界、改 value 语义，还是换算法。

\noindent\textbf{LeetCode 46 全排列。}

\textbf{题目说明。} LeetCode 46 全排列 的题面入口要先还原成具体任务：排列关心元素顺序，每个位置从所有未使用元素中选择一个。

\textbf{样例入口。} 拿 [1, 2, 3] 手算，第一位可选 1、2、3；选 1 后，第二位只能从未使用的 2、3 中选。

\textbf{暴力基线。} 暴力可以枚举所有长度为 n 的序列，再检查是否每个元素刚好出现一次。

\textbf{暴力过程。} 以 [1, 2, 3] 为例，第一位可选 1、2、3；若第一位选 1，第二位只能从 2、3 中选，第三位由剩余元素决定。

\textbf{重复工作。} 题面模型：排列关心元素顺序，每个位置从所有未使用元素中选择一个。暴力版的慢点要落到本题：浪费在生成了含重复使用元素的非法序列。排列的每一层只需要从未使用元素中选择。后面的优化只消掉这类重复工作，不能改变答案集合。

\textbf{优化条件。} 本题的关键观察是：路径长度达到 n 时产生答案，used 数组保证同一元素不重复使用。这句话必须能翻译成可维护状态；如果题目条件改变后观察不再成立，就不能继续套原来的写法。

\textbf{相邻状态。} 规律是路径长度等于 n 时得到一个排列，used 数组表示哪些下标已经在当前路径里。把这个特例继续拆开看：一层递归对应一个明确决策，路径保存已经做出的选择，选择列表保存仍可能扩展的分支。剪枝前必须能说明被剪掉的整棵子树没有合法答案，而不是只因为它看起来不优。

\textbf{状态复用。} 回溯路径 path 保存当前排列前缀，used[i] 表示 nums[i] 是否已在路径中。path 长度等于 n 时复制进答案。手算时只改被新输入影响的那一小部分，而不是回头重算完整候选。

\textbf{指针和字段。} 状态字段要能说清：depth 是当前要填的位置，path 是已选择序列，used 是可选状态，answers 保存所有完整排列。手算时按这个协议跟踪：[1, 2, 3] 中路径 [1] 后，used[1]=true；下一层只能选 2 或 3。回退时弹出 1 并恢复 used。

\textbf{代码落点。} 关键是进入分支前 path.push\_back(nums[i]); used[i]=true；退出分支时反向恢复。保存答案时必须复制 path。关键代码行必须能逐句翻译成“旧状态怎样变成新状态”，否则就只是模板。

\textbf{正确性。} 每个位置都枚举所有未使用元素，所以任意合法排列都有唯一生成路径；used 阻止同一元素在一条路径中重复出现。

\textbf{边界实验。} 先用空数组、单元素、输出规模为 n 阶乘、元素值唯一假设做反例检查。实验时覆盖空数组、单元素、输出规模为 n 阶乘、元素值唯一假设，统计搜索节点数量和剪枝次数。重复元素题要打印同一层跳过哪些值，确认没有误删合法路径。

\textbf{复杂度变化。} 输出规模为 n! 个排列，每个排列长度 n，时间 O(n*n!)，空间 O(n) 不计输出。

\textbf{迁移追问。} 如果题目改成“若有重复元素，需要排序并在同层跳过等价选择”，先判断原状态是否仍然覆盖新答案集合，再决定是微调边界、改 value 语义，还是换算法。

\noindent\textbf{LeetCode 47 全排列 II。}

\textbf{题目说明。} LeetCode 47 全排列 II 的题面入口要先还原成具体任务：重复元素会让不同下标选择生成相同排列。

\textbf{样例入口。} 拿 [1, 1, 2] 手算，第一位选第一个 1 和第二个 1 会生成同样分支，所以同层必须跳过后一个 1；但第一层选 1 后，第二层仍可选另一个 1。

\textbf{暴力基线。} 慢版本就是完整搜索树：每层列出选择列表，路径到达终止条件时检查答案。它先保证覆盖所有可能，再把明显无效或重复的分支剪掉。放到本题就是：先按“重复元素会让不同下标选择生成相同排列”完整试慢版本；样例里已经能看到“规律是排序后用 used 区分路径内使用，同层重复值只保留第一个未使用代表”，所以优化前必须先能说清慢版本枚举了哪些候选。

\textbf{暴力过程。} 拿 [1, 1, 2] 手算，第一位选第一个 1 和第二个 1 会生成同样分支，所以同层必须跳过后一个 1；但第一层选 1 后，第二层仍可选另一个 1。

\textbf{重复工作。} 题面模型：重复元素会让不同下标选择生成相同排列。暴力版的慢点要落到本题：慢版本会围绕“重复元素会让不同下标选择生成相同排列”展开完整搜索树。样例规律是“规律是排序后用 used 区分路径内使用，同层重复值只保留第一个未使用代表”，说明一层递归必须对应一个明确决策，剪枝只能删除整棵不可能成功或重复的子树。后面的优化只消掉这类重复工作，不能改变答案集合。

\textbf{优化条件。} 本题的关键观察是：排序后使用 used 数组，同层遇到相同值且前一个相同值未使用时跳过。这句话必须能翻译成可维护状态；如果题目条件改变后观察不再成立，就不能继续套原来的写法。

\textbf{相邻状态。} 规律是排序后用 used 区分路径内使用，同层重复值只保留第一个未使用代表。把这个特例继续拆开看：一层递归对应一个明确决策，路径保存已经做出的选择，选择列表保存仍可能扩展的分支。剪枝前必须能说明被剪掉的整棵子树没有合法答案，而不是只因为它看起来不优。

\textbf{状态复用。} 把完整搜索树加上合法性检查、剩余资源剪枝和同层去重；路径状态进入和退出要成对恢复。本题落点是：规律是排序后用 used 区分路径内使用，同层重复值只保留第一个未使用代表。手算时只改被新输入影响的那一小部分，而不是回头重算完整候选。

\textbf{指针和字段。} 状态字段要能说清：路径、起点或使用标记、剩余目标、答案集合、剪枝条件；进入递归和退出递归必须成对恢复。手算时按这个协议跟踪：拿 [1, 1, 2] 手算，第一位选第一个 1 和第二个 1 会生成同样分支，所以同层必须跳过后一个 1；但第一层选 1 后，第二层仍可选另一个 1。然后按协议复查：手算时给每层标出选择列表和路径。剪枝发生前要指出被剪分支为什么已经不可能得到合法答案。

\textbf{代码落点。} 剪枝条件写在扩展前；同层去重不要误写成全局去重；保存答案时复制当前路径。关键代码行必须能逐句翻译成“旧状态怎样变成新状态”，否则就只是模板。

\textbf{正确性。} 证明从搜索树覆盖开始：每个合法答案有且只有一条路径，剪枝只删掉不可能成功的分支。

\textbf{边界实验。} 先用全重复、部分重复、同层去重条件写反、输出顺序不要求做反例检查。实验时覆盖全重复、部分重复、同层去重条件写反、输出顺序不要求，统计搜索节点数量和剪枝次数。重复元素题要打印同一层跳过哪些值，确认没有误删合法路径。

\textbf{复杂度变化。} 暴力搜索树的规模由输出或选择空间决定；剪枝不改变最坏指数级本质，但能删除不可能成功或重复的分支，复杂度要说明输出规模。

\textbf{迁移追问。} 如果题目改成“子集 II 和组合总和 II 也是同层去重，但起点推进方式不同”，先判断原状态是否仍然覆盖新答案集合，再决定是微调边界、改 value 语义，还是换算法。

\noindent\textbf{LeetCode 51 N 皇后。}

\textbf{题目说明。} LeetCode 51 N 皇后 的题面入口要先还原成具体任务：每一行放一个皇后，列和两条对角线不能冲突。

\textbf{样例入口。} 拿 n=4 手算，第一行放某列后，下一行不能放同列、主对角线和副对角线冲突的位置。

\textbf{暴力基线。} 慢版本就是完整搜索树：每层列出选择列表，路径到达终止条件时检查答案。它先保证覆盖所有可能，再把明显无效或重复的分支剪掉。放到本题就是：先按“每一行放一个皇后，列和两条对角线不能冲突”完整试慢版本；样例里已经能看到“规律是按行递归，每行放一个皇后，三个集合分别记录列、row-col、row+col 是否被占用”，所以优化前必须先能说清慢版本枚举了哪些候选。

\textbf{暴力过程。} 拿 n=4 手算，第一行放某列后，下一行不能放同列、主对角线和副对角线冲突的位置。

\textbf{重复工作。} 题面模型：每一行放一个皇后，列和两条对角线不能冲突。暴力版的慢点要落到本题：慢版本会围绕“每一行放一个皇后，列和两条对角线不能冲突”展开完整搜索树。样例规律是“规律是按行递归，每行放一个皇后，三个集合分别记录列、row-col、row+col 是否被占用”，说明一层递归必须对应一个明确决策，剪枝只能删除整棵不可能成功或重复的子树。后面的优化只消掉这类重复工作，不能改变答案集合。

\textbf{优化条件。} 本题的关键观察是：逐行搜索，用列集合、主对角线集合和副对角线集合快速判断合法性。这句话必须能翻译成可维护状态；如果题目条件改变后观察不再成立，就不能继续套原来的写法。

\textbf{相邻状态。} 规律是按行递归，每行放一个皇后，三个集合分别记录列、row-col、row+col 是否被占用。把这个特例继续拆开看：一层递归对应一个明确决策，路径保存已经做出的选择，选择列表保存仍可能扩展的分支。剪枝前必须能说明被剪掉的整棵子树没有合法答案，而不是只因为它看起来不优。

\textbf{状态复用。} 把完整搜索树加上合法性检查、剩余资源剪枝和同层去重；路径状态进入和退出要成对恢复。本题落点是：规律是按行递归，每行放一个皇后，三个集合分别记录列、row-col、row+col 是否被占用。手算时只改被新输入影响的那一小部分，而不是回头重算完整候选。

\textbf{指针和字段。} 状态字段要能说清：路径、起点或使用标记、剩余目标、答案集合、剪枝条件；进入递归和退出递归必须成对恢复。手算时按这个协议跟踪：拿 n=4 手算，第一行放某列后，下一行不能放同列、主对角线和副对角线冲突的位置。然后按协议复查：手算时给每层标出选择列表和路径。剪枝发生前要指出被剪分支为什么已经不可能得到合法答案。

\textbf{代码落点。} 剪枝条件写在扩展前；同层去重不要误写成全局去重；保存答案时复制当前路径。关键代码行必须能逐句翻译成“旧状态怎样变成新状态”，否则就只是模板。

\textbf{正确性。} 证明从搜索树覆盖开始：每个合法答案有且只有一条路径，剪枝只删掉不可能成功的分支。

\textbf{边界实验。} 先用n 为一、n 为二或三无解、对角线编号、回退时状态恢复做反例检查。实验时覆盖n 为一、n 为二或三无解、对角线编号、回退时状态恢复，统计搜索节点数量和剪枝次数。重复元素题要打印同一层跳过哪些值，确认没有误删合法路径。

\textbf{复杂度变化。} 暴力搜索树的规模由输出或选择空间决定；剪枝不改变最坏指数级本质，但能删除不可能成功或重复的分支，复杂度要说明输出规模。

\textbf{迁移追问。} 如果题目改成“可用位掩码压缩三个约束集合，提高常数性能”，先判断原状态是否仍然覆盖新答案集合，再决定是微调边界、改 value 语义，还是换算法。

\noindent\textbf{LeetCode 78 子集。}

\textbf{题目说明。} LeetCode 78 子集 的题面入口要先还原成具体任务：每个元素都有选或不选两种选择，答案是所有路径节点。

\textbf{样例入口。} 拿 [1, 2] 画搜索树，第一层可以不选 1 或选 1；在选 1 的分支里再决定 2，得到 [1] 和 [1, 2]。

\textbf{暴力基线。} 暴力可以枚举长度从 0 到 n 的所有组合，并为每种长度单独搜索。

\textbf{暴力过程。} 以 [1, 2, 3] 为例，空集、[1]、[2]、[3]、[1, 2] 等都是答案；一个元素只有选和不选两种状态。

\textbf{重复工作。} 题面模型：每个元素都有选或不选两种选择，答案是所有路径节点。暴力版的慢点要落到本题：浪费在按长度重复组织搜索。子集问题的每个搜索树节点本身就是一个答案，不必等到固定长度。后面的优化只消掉这类重复工作，不能改变答案集合。

\textbf{优化条件。} 本题的关键观察是：回溯按起点向后枚举，避免不同顺序生成同一集合。这句话必须能翻译成可维护状态；如果题目条件改变后观察不再成立，就不能继续套原来的写法。

\textbf{相邻状态。} 规律是每个元素对应选或不选，或用起点向后枚举，路径上的每个节点都是一个子集；空数组也要输出空集。把这个特例继续拆开看：一层递归对应一个明确决策，路径保存已经做出的选择，选择列表保存仍可能扩展的分支。剪枝前必须能说明被剪掉的整棵子树没有合法答案，而不是只因为它看起来不优。

\textbf{状态复用。} 回溯时先把当前 path 加入答案，再从 start 到末尾选择下一个元素，递归时 start=i+1，保证元素顺序递增。手算时只改被新输入影响的那一小部分，而不是回头重算完整候选。

\textbf{指针和字段。} 状态字段要能说清：start 是下一次可选择的最小下标，path 是当前子集，answers 保存所有已访问路径节点。手算时按这个协议跟踪：从空集开始记录 []；选择 1 后记录 [1]；再选择 2 记录 [1, 2]；回退后选择 3 得到 [1, 3]。

\textbf{代码落点。} 关键是每进入一个节点就记录 path，而不是只在叶子记录。递归下一层用 i+1，避免 [1, 2] 和 [2, 1] 重复。关键代码行必须能逐句翻译成“旧状态怎样变成新状态”，否则就只是模板。

\textbf{正确性。} 任意子集按原数组下标升序排列后，对应搜索树中唯一一条递增路径；每个路径节点都被记录，所以不漏不重。

\textbf{边界实验。} 先用空数组、单元素、输出数量为二的 n 次方、顺序不要求做反例检查。实验时覆盖空数组、单元素、输出数量为二的 n 次方、顺序不要求，统计搜索节点数量和剪枝次数。重复元素题要打印同一层跳过哪些值，确认没有误删合法路径。

\textbf{复杂度变化。} 共有 \ensuremath{2^{n}} 个子集，复制路径总成本 O(\ensuremath{n \cdot 2^{n}})，递归栈 O(n)，输出空间另计。

\textbf{迁移追问。} 如果题目改成“也可用位掩码枚举所有子集，适合 n 较小的场景”，先判断原状态是否仍然覆盖新答案集合，再决定是微调边界、改 value 语义，还是换算法。

\noindent\textbf{LeetCode 79 单词搜索。}

\textbf{题目说明。} LeetCode 79 单词搜索 的题面入口要先还原成具体任务：网格路径不能重复使用同一格，状态包含位置、匹配下标和访问标记。

\textbf{样例入口。} 拿 board=[A B; C D]、word=AB 手算，从 A 出发只能走到相邻的 B，不能跳到 D；若搜索 ABA，第二个 A 不能复用起点格。

\textbf{暴力基线。} 慢版本就是完整搜索树：每层列出选择列表，路径到达终止条件时检查答案。它先保证覆盖所有可能，再把明显无效或重复的分支剪掉。放到本题就是：先按“网格路径不能重复使用同一格，状态包含位置、匹配下标和访问标记”完整试慢版本；样例里已经能看到“规律是状态包含位置、匹配下标和访问标记，递归退出时要恢复访问标记”，所以优化前必须先能说清慢版本枚举了哪些候选。

\textbf{暴力过程。} 拿 board=[A B; C D]、word=AB 手算，从 A 出发只能走到相邻的 B，不能跳到 D；若搜索 ABA，第二个 A 不能复用起点格。

\textbf{重复工作。} 题面模型：网格路径不能重复使用同一格，状态包含位置、匹配下标和访问标记。暴力版的慢点要落到本题：慢版本会围绕“网格路径不能重复使用同一格，状态包含位置、匹配下标和访问标记”展开完整搜索树。样例规律是“规律是状态包含位置、匹配下标和访问标记，递归退出时要恢复访问标记”，说明一层递归必须对应一个明确决策，剪枝只能删除整棵不可能成功或重复的子树。后面的优化只消掉这类重复工作，不能改变答案集合。

\textbf{优化条件。} 本题的关键观察是：剪枝包括首字符不匹配、剩余长度不足、字符频次不足和越界访问。这句话必须能翻译成可维护状态；如果题目条件改变后观察不再成立，就不能继续套原来的写法。

\textbf{相邻状态。} 规律是状态包含位置、匹配下标和访问标记，递归退出时要恢复访问标记。把这个特例继续拆开看：一层递归对应一个明确决策，路径保存已经做出的选择，选择列表保存仍可能扩展的分支。剪枝前必须能说明被剪掉的整棵子树没有合法答案，而不是只因为它看起来不优。

\textbf{状态复用。} 把完整搜索树加上合法性检查、剩余资源剪枝和同层去重；路径状态进入和退出要成对恢复。本题落点是：规律是状态包含位置、匹配下标和访问标记，递归退出时要恢复访问标记。手算时只改被新输入影响的那一小部分，而不是回头重算完整候选。

\textbf{指针和字段。} 状态字段要能说清：路径、起点或使用标记、剩余目标、答案集合、剪枝条件；进入递归和退出递归必须成对恢复。手算时按这个协议跟踪：拿 board=[A B; C D]、word=AB 手算，从 A 出发只能走到相邻的 B，不能跳到 D；若搜索 ABA，第二个 A 不能复用起点格。然后按协议复查：手算时给每层标出选择列表和路径。剪枝发生前要指出被剪分支为什么已经不可能得到合法答案。

\textbf{代码落点。} 剪枝条件写在扩展前；同层去重不要误写成全局去重；保存答案时复制当前路径。关键代码行必须能逐句翻译成“旧状态怎样变成新状态”，否则就只是模板。

\textbf{正确性。} 证明从搜索树覆盖开始：每个合法答案有且只有一条路径，剪枝只删掉不可能成功的分支。

\textbf{边界实验。} 先用单字符单词、路径需要回退、重复字母、单元格不能复用做反例检查。实验时覆盖单字符单词、路径需要回退、重复字母、单元格不能复用，统计搜索节点数量和剪枝次数。重复元素题要打印同一层跳过哪些值，确认没有误删合法路径。

\textbf{复杂度变化。} 暴力搜索树的规模由输出或选择空间决定；剪枝不改变最坏指数级本质，但能删除不可能成功或重复的分支，复杂度要说明输出规模。

\textbf{迁移追问。} 如果题目改成“若要查很多单词，应使用 Trie 合并搜索前缀”，先判断原状态是否仍然覆盖新答案集合，再决定是微调边界、改 value 语义，还是换算法。

\noindent\textbf{LeetCode 90 子集 II。}

\textbf{题目说明。} LeetCode 90 子集 II 的题面入口要先还原成具体任务：重复元素会让不同下标路径生成相同集合。

\textbf{样例入口。} 拿 [1, 2, 2] 手算，排序后第二个 2 如果在同一层再次作为起点，会生成和第一个 2 同样的子集；但在已经选择第一个 2 的下一层，第二个 2 又可以被选出 [2, 2]。

\textbf{暴力基线。} 慢版本就是完整搜索树：每层列出选择列表，路径到达终止条件时检查答案。它先保证覆盖所有可能，再把明显无效或重复的分支剪掉。放到本题就是：先按“重复元素会让不同下标路径生成相同集合”完整试慢版本；样例里已经能看到“规律是同层跳过重复值，不是全局禁止重复值”，所以优化前必须先能说清慢版本枚举了哪些候选。

\textbf{暴力过程。} 拿 [1, 2, 2] 手算，排序后第二个 2 如果在同一层再次作为起点，会生成和第一个 2 同样的子集；但在已经选择第一个 2 的下一层，第二个 2 又可以被选出 [2, 2]。

\textbf{重复工作。} 题面模型：重复元素会让不同下标路径生成相同集合。暴力版的慢点要落到本题：慢版本会围绕“重复元素会让不同下标路径生成相同集合”展开完整搜索树。样例规律是“规律是同层跳过重复值，不是全局禁止重复值”，说明一层递归必须对应一个明确决策，剪枝只能删除整棵不可能成功或重复的子树。后面的优化只消掉这类重复工作，不能改变答案集合。

\textbf{优化条件。} 本题的关键观察是：排序后在同一层跳过相同值，保留不同层使用重复值的可能。这句话必须能翻译成可维护状态；如果题目条件改变后观察不再成立，就不能继续套原来的写法。

\textbf{相邻状态。} 规律是同层跳过重复值，不是全局禁止重复值。把这个特例继续拆开看：一层递归对应一个明确决策，路径保存已经做出的选择，选择列表保存仍可能扩展的分支。剪枝前必须能说明被剪掉的整棵子树没有合法答案，而不是只因为它看起来不优。

\textbf{状态复用。} 把完整搜索树加上合法性检查、剩余资源剪枝和同层去重；路径状态进入和退出要成对恢复。本题落点是：规律是同层跳过重复值，不是全局禁止重复值。手算时只改被新输入影响的那一小部分，而不是回头重算完整候选。

\textbf{指针和字段。} 状态字段要能说清：路径、起点或使用标记、剩余目标、答案集合、剪枝条件；进入递归和退出递归必须成对恢复。手算时按这个协议跟踪：拿 [1, 2, 2] 手算，排序后第二个 2 如果在同一层再次作为起点，会生成和第一个 2 同样的子集；但在已经选择第一个 2 的下一层，第二个 2 又可以被选出 [2, 2]。然后按协议复查：手算时给每层标出选择列表和路径。剪枝发生前要指出被剪分支为什么已经不可能得到合法答案。

\textbf{代码落点。} 剪枝条件写在扩展前；同层去重不要误写成全局去重；保存答案时复制当前路径。关键代码行必须能逐句翻译成“旧状态怎样变成新状态”，否则就只是模板。

\textbf{正确性。} 证明从搜索树覆盖开始：每个合法答案有且只有一条路径，剪枝只删掉不可能成功的分支。

\textbf{边界实验。} 先用全重复、空集、重复值相邻、去重条件写成同层而不是全局做反例检查。实验时覆盖全重复、空集、重复值相邻、去重条件写成同层而不是全局，统计搜索节点数量和剪枝次数。重复元素题要打印同一层跳过哪些值，确认没有误删合法路径。

\textbf{复杂度变化。} 暴力搜索树的规模由输出或选择空间决定；剪枝不改变最坏指数级本质，但能删除不可能成功或重复的分支，复杂度要说明输出规模。

\textbf{迁移追问。} 如果题目改成“组合总和 II 使用同样的同层去重思想”，先判断原状态是否仍然覆盖新答案集合，再决定是微调边界、改 value 语义，还是换算法。

\noindent\textbf{LeetCode 93 复原 IP 地址。}

\textbf{题目说明。} LeetCode 93 复原 IP 地址 的题面入口要先还原成具体任务：字符串要切成四段，每段必须是 0 到 255 且无非法前导零。

\textbf{样例入口。} 拿 25525511135 手算，第一段可以是 255，但不能是 2552；段值必须在 0..255，且 01 这种前导零非法。

\textbf{暴力基线。} 慢版本就是完整搜索树：每层列出选择列表，路径到达终止条件时检查答案。它先保证覆盖所有可能，再把明显无效或重复的分支剪掉。放到本题就是：先按“字符串要切成四段，每段必须是 0 到 255 且无非法前导零”完整试慢版本；样例里已经能看到“规律是回溯选四段，每段长度 1 到 3，剩余字符数必须能填满剩余段数”，所以优化前必须先能说清慢版本枚举了哪些候选。

\textbf{暴力过程。} 拿 25525511135 手算，第一段可以是 255，但不能是 2552；段值必须在 0..255，且 01 这种前导零非法。

\textbf{重复工作。} 题面模型：字符串要切成四段，每段必须是 0 到 255 且无非法前导零。暴力版的慢点要落到本题：慢版本会围绕“字符串要切成四段，每段必须是 0 到 255 且无非法前导零”展开完整搜索树。样例规律是“规律是回溯选四段，每段长度 1 到 3，剩余字符数必须能填满剩余段数”，说明一层递归必须对应一个明确决策，剪枝只能删除整棵不可能成功或重复的子树。后面的优化只消掉这类重复工作，不能改变答案集合。

\textbf{优化条件。} 本题的关键观察是：按段数和剩余长度剪枝，每次尝试长度一到三的片段。这句话必须能翻译成可维护状态；如果题目条件改变后观察不再成立，就不能继续套原来的写法。

\textbf{相邻状态。} 规律是回溯选四段，每段长度 1 到 3，剩余字符数必须能填满剩余段数。把这个特例继续拆开看：一层递归对应一个明确决策，路径保存已经做出的选择，选择列表保存仍可能扩展的分支。剪枝前必须能说明被剪掉的整棵子树没有合法答案，而不是只因为它看起来不优。

\textbf{状态复用。} 把完整搜索树加上合法性检查、剩余资源剪枝和同层去重；路径状态进入和退出要成对恢复。本题落点是：规律是回溯选四段，每段长度 1 到 3，剩余字符数必须能填满剩余段数。手算时只改被新输入影响的那一小部分，而不是回头重算完整候选。

\textbf{指针和字段。} 状态字段要能说清：路径、起点或使用标记、剩余目标、答案集合、剪枝条件；进入递归和退出递归必须成对恢复。手算时按这个协议跟踪：拿 25525511135 手算，第一段可以是 255，但不能是 2552；段值必须在 0..255，且 01 这种前导零非法。然后按协议复查：手算时给每层标出选择列表和路径。剪枝发生前要指出被剪分支为什么已经不可能得到合法答案。

\textbf{代码落点。} 剪枝条件写在扩展前；同层去重不要误写成全局去重；保存答案时复制当前路径。关键代码行必须能逐句翻译成“旧状态怎样变成新状态”，否则就只是模板。

\textbf{正确性。} 证明从搜索树覆盖开始：每个合法答案有且只有一条路径，剪枝只删掉不可能成功的分支。

\textbf{边界实验。} 先用前导零、数值超过 255、剩余字符太多或太少、刚好四段做反例检查。实验时覆盖前导零、数值超过 255、剩余字符太多或太少、刚好四段，统计搜索节点数量和剪枝次数。重复元素题要打印同一层跳过哪些值，确认没有误删合法路径。

\textbf{复杂度变化。} 暴力搜索树的规模由输出或选择空间决定；剪枝不改变最坏指数级本质，但能删除不可能成功或重复的分支，复杂度要说明输出规模。

\textbf{迁移追问。} 如果题目改成“分割回文串也是前缀分割，但合法性检查不同”，先判断原状态是否仍然覆盖新答案集合，再决定是微调边界、改 value 语义，还是换算法。

\noindent\textbf{LeetCode 131 分割回文串。}

\textbf{题目说明。} LeetCode 131 分割回文串 的题面入口要先还原成具体任务：每次选择一个从当前位置开始的回文前缀。

\textbf{样例入口。} 拿 aab 手算，第一刀可以切 a，后面 ab 不全是回文；也可以切 aa，后面 b 是回文，得到 [aa, b]。

\textbf{暴力基线。} 慢版本就是完整搜索树：每层列出选择列表，路径到达终止条件时检查答案。它先保证覆盖所有可能，再把明显无效或重复的分支剪掉。放到本题就是：先按“每次选择一个从当前位置开始的回文前缀”完整试慢版本；样例里已经能看到“规律是回溯枚举切分点，但每段必须先通过回文检查；预处理回文表可以避免重复检查子串”，所以优化前必须先能说清慢版本枚举了哪些候选。

\textbf{暴力过程。} 拿 aab 手算，第一刀可以切 a，后面 ab 不全是回文；也可以切 aa，后面 b 是回文，得到 [aa, b]。

\textbf{重复工作。} 题面模型：每次选择一个从当前位置开始的回文前缀。暴力版的慢点要落到本题：慢版本会围绕“每次选择一个从当前位置开始的回文前缀”展开完整搜索树。样例规律是“规律是回溯枚举切分点，但每段必须先通过回文检查；预处理回文表可以避免重复检查子串”，说明一层递归必须对应一个明确决策，剪枝只能删除整棵不可能成功或重复的子树。后面的优化只消掉这类重复工作，不能改变答案集合。

\textbf{优化条件。} 本题的关键观察是：预处理回文 DP 表，让回溯合法性检查为常数。这句话必须能翻译成可维护状态；如果题目条件改变后观察不再成立，就不能继续套原来的写法。

\textbf{相邻状态。} 规律是回溯枚举切分点，但每段必须先通过回文检查；预处理回文表可以避免重复检查子串。把这个特例继续拆开看：一层递归对应一个明确决策，路径保存已经做出的选择，选择列表保存仍可能扩展的分支。剪枝前必须能说明被剪掉的整棵子树没有合法答案，而不是只因为它看起来不优。

\textbf{状态复用。} 把完整搜索树加上合法性检查、剩余资源剪枝和同层去重；路径状态进入和退出要成对恢复。本题落点是：规律是回溯枚举切分点，但每段必须先通过回文检查；预处理回文表可以避免重复检查子串。手算时只改被新输入影响的那一小部分，而不是回头重算完整候选。

\textbf{指针和字段。} 状态字段要能说清：路径、起点或使用标记、剩余目标、答案集合、剪枝条件；进入递归和退出递归必须成对恢复。手算时按这个协议跟踪：拿 aab 手算，第一刀可以切 a，后面 ab 不全是回文；也可以切 aa，后面 b 是回文，得到 [aa, b]。然后按协议复查：手算时给每层标出选择列表和路径。剪枝发生前要指出被剪分支为什么已经不可能得到合法答案。

\textbf{代码落点。} 剪枝条件写在扩展前；同层去重不要误写成全局去重；保存答案时复制当前路径。关键代码行必须能逐句翻译成“旧状态怎样变成新状态”，否则就只是模板。

\textbf{正确性。} 证明从搜索树覆盖开始：每个合法答案有且只有一条路径，剪枝只删掉不可能成功的分支。

\textbf{边界实验。} 先用空串、单字符、全相同字符、重复分割路径做反例检查。实验时覆盖空串、单字符、全相同字符、重复分割路径，统计搜索节点数量和剪枝次数。重复元素题要打印同一层跳过哪些值，确认没有误删合法路径。

\textbf{复杂度变化。} 暴力搜索树的规模由输出或选择空间决定；剪枝不改变最坏指数级本质，但能删除不可能成功或重复的分支，复杂度要说明输出规模。

\textbf{迁移追问。} 如果题目改成“若只求最少切割次数，可改成 DP 最短切分”，先判断原状态是否仍然覆盖新答案集合，再决定是微调边界、改 value 语义，还是换算法。

\noindent\textbf{LeetCode 140 单词拆分 II。}

\textbf{题目说明。} LeetCode 140 单词拆分 II 的题面入口要先还原成具体任务：需要输出所有由字典单词拼出的句子，输出规模本身可能指数级。

\textbf{样例入口。} 拿 catsanddog 手算，前缀 cat 和 cats 都能走，但某些后缀继续切会失败。若直接 DFS，每次都会重复探索同一个后缀。

\textbf{暴力基线。} 慢版本就是完整搜索树：每层列出选择列表，路径到达终止条件时检查答案。它先保证覆盖所有可能，再把明显无效或重复的分支剪掉。放到本题就是：先按“需要输出所有由字典单词拼出的句子，输出规模本身可能指数级”完整试慢版本；样例里已经能看到“规律是先用记忆化记录某个起点之后能生成哪些句子，或先用 DP 剪掉不可拆后缀，再枚举路径”，所以优化前必须先能说清慢版本枚举了哪些候选。

\textbf{暴力过程。} 拿 catsanddog 手算，前缀 cat 和 cats 都能走，但某些后缀继续切会失败。若直接 DFS，每次都会重复探索同一个后缀。

\textbf{重复工作。} 题面模型：需要输出所有由字典单词拼出的句子，输出规模本身可能指数级。暴力版的慢点要落到本题：慢版本会围绕“需要输出所有由字典单词拼出的句子，输出规模本身可能指数级”展开完整搜索树。样例规律是“规律是先用记忆化记录某个起点之后能生成哪些句子，或先用 DP 剪掉不可拆后缀，再枚举路径”，说明一层递归必须对应一个明确决策，剪枝只能删除整棵不可能成功或重复的子树。后面的优化只消掉这类重复工作，不能改变答案集合。

\textbf{优化条件。} 本题的关键观察是：先用 DP 或记忆化判断后缀是否可拆，再 DFS 枚举可行前缀并在末尾拼接句子。这句话必须能翻译成可维护状态；如果题目条件改变后观察不再成立，就不能继续套原来的写法。

\textbf{相邻状态。} 规律是先用记忆化记录某个起点之后能生成哪些句子，或先用 DP 剪掉不可拆后缀，再枚举路径。把这个特例继续拆开看：一层递归对应一个明确决策，路径保存已经做出的选择，选择列表保存仍可能扩展的分支。剪枝前必须能说明被剪掉的整棵子树没有合法答案，而不是只因为它看起来不优。

\textbf{状态复用。} 把完整搜索树加上合法性检查、剩余资源剪枝和同层去重；路径状态进入和退出要成对恢复。本题落点是：规律是先用记忆化记录某个起点之后能生成哪些句子，或先用 DP 剪掉不可拆后缀，再枚举路径。手算时只改被新输入影响的那一小部分，而不是回头重算完整候选。

\textbf{指针和字段。} 状态字段要能说清：路径、起点或使用标记、剩余目标、答案集合、剪枝条件；进入递归和退出递归必须成对恢复。手算时按这个协议跟踪：拿 catsanddog 手算，前缀 cat 和 cats 都能走，但某些后缀继续切会失败。若直接 DFS，每次都会重复探索同一个后缀。然后按协议复查：手算时给每层标出选择列表和路径。剪枝发生前要指出被剪分支为什么已经不可能得到合法答案。

\textbf{代码落点。} 剪枝条件写在扩展前；同层去重不要误写成全局去重；保存答案时复制当前路径。关键代码行必须能逐句翻译成“旧状态怎样变成新状态”，否则就只是模板。

\textbf{正确性。} 证明从搜索树覆盖开始：每个合法答案有且只有一条路径，剪枝只删掉不可能成功的分支。

\textbf{边界实验。} 先用无解后缀、重复单词、路径拼接成本、答案数量主导复杂度做反例检查。实验时覆盖无解后缀、重复单词、路径拼接成本、答案数量主导复杂度，统计搜索节点数量和剪枝次数。重复元素题要打印同一层跳过哪些值，确认没有误删合法路径。

\textbf{复杂度变化。} 暴力搜索树的规模由输出或选择空间决定；剪枝不改变最坏指数级本质，但能删除不可能成功或重复的分支，复杂度要说明输出规模。

\textbf{迁移追问。} 如果题目改成“只问能否拆分时用布尔 DP 即可，不要承担枚举所有句子的成本”，先判断原状态是否仍然覆盖新答案集合，再决定是微调边界、改 value 语义，还是换算法。

\noindent\textbf{LeetCode 212 单词搜索 II。}

\textbf{题目说明。} LeetCode 212 单词搜索 II 的题面入口要先还原成具体任务：多个单词在同一棋盘上搜索，公共前缀可以共享。

\textbf{样例入口。} 拿 board 上有 oath、pea、eat 手算，单词共享前缀时不应对每个词重复全图搜索；从 o 出发沿 Trie 前缀 o-a-t-h 找到 oath。

\textbf{暴力基线。} 慢版本就是完整搜索树：每层列出选择列表，路径到达终止条件时检查答案。它先保证覆盖所有可能，再把明显无效或重复的分支剪掉。放到本题就是：先按“多个单词在同一棋盘上搜索，公共前缀可以共享”完整试慢版本；样例里已经能看到“规律是 Trie 合并所有单词前缀，DFS 走棋盘同时走 Trie，命中单词后记录并可去重”，所以优化前必须先能说清慢版本枚举了哪些候选。

\textbf{暴力过程。} 拿 board 上有 oath、pea、eat 手算，单词共享前缀时不应对每个词重复全图搜索；从 o 出发沿 Trie 前缀 o-a-t-h 找到 oath。

\textbf{重复工作。} 题面模型：多个单词在同一棋盘上搜索，公共前缀可以共享。暴力版的慢点要落到本题：慢版本会围绕“多个单词在同一棋盘上搜索，公共前缀可以共享”展开完整搜索树。样例规律是“规律是 Trie 合并所有单词前缀，DFS 走棋盘同时走 Trie，命中单词后记录并可去重”，说明一层递归必须对应一个明确决策，剪枝只能删除整棵不可能成功或重复的子树。后面的优化只消掉这类重复工作，不能改变答案集合。

\textbf{优化条件。} 本题的关键观察是：先建 Trie，再从每个格子回溯；路径到达单词结束节点就收集答案。这句话必须能翻译成可维护状态；如果题目条件改变后观察不再成立，就不能继续套原来的写法。

\textbf{相邻状态。} 规律是 Trie 合并所有单词前缀，DFS 走棋盘同时走 Trie，命中单词后记录并可去重。把这个特例继续拆开看：一层递归对应一个明确决策，路径保存已经做出的选择，选择列表保存仍可能扩展的分支。剪枝前必须能说明被剪掉的整棵子树没有合法答案，而不是只因为它看起来不优。

\textbf{状态复用。} 把完整搜索树加上合法性检查、剩余资源剪枝和同层去重；路径状态进入和退出要成对恢复。本题落点是：规律是 Trie 合并所有单词前缀，DFS 走棋盘同时走 Trie，命中单词后记录并可去重。手算时只改被新输入影响的那一小部分，而不是回头重算完整候选。

\textbf{指针和字段。} 状态字段要能说清：路径、起点或使用标记、剩余目标、答案集合、剪枝条件；进入递归和退出递归必须成对恢复。手算时按这个协议跟踪：拿 board 上有 oath、pea、eat 手算，单词共享前缀时不应对每个词重复全图搜索；从 o 出发沿 Trie 前缀 o-a-t-h 找到 oath。然后按协议复查：手算时给每层标出选择列表和路径。剪枝发生前要指出被剪分支为什么已经不可能得到合法答案。

\textbf{代码落点。} 剪枝条件写在扩展前；同层去重不要误写成全局去重；保存答案时复制当前路径。关键代码行必须能逐句翻译成“旧状态怎样变成新状态”，否则就只是模板。

\textbf{正确性。} 证明从搜索树覆盖开始：每个合法答案有且只有一条路径，剪枝只删掉不可能成功的分支。

\textbf{边界实验。} 先用重复单词、同一格不能复用、前缀是完整单词、剪枝删除空子树做反例检查。实验时覆盖重复单词、同一格不能复用、前缀是完整单词、剪枝删除空子树，统计搜索节点数量和剪枝次数。重复元素题要打印同一层跳过哪些值，确认没有误删合法路径。

\textbf{复杂度变化。} 暴力搜索树的规模由输出或选择空间决定；剪枝不改变最坏指数级本质，但能删除不可能成功或重复的分支，复杂度要说明输出规模。

\textbf{迁移追问。} 如果题目改成“这是单词搜索从单目标扩展到多目标后的结构升级”，先判断原状态是否仍然覆盖新答案集合，再决定是微调边界、改 value 语义，还是换算法。

\noindent\textbf{LeetCode 216 组合总和 III。}

\textbf{题目说明。} LeetCode 216 组合总和 III 的题面入口要先还原成具体任务：从一到九中选 k 个数和为 n，每个数最多用一次。

\textbf{样例入口。} 拿 k=3、n=7 手算，选择 1 后还需两个数凑 6，可以选 2 和 4；若当前和已经超过 7，就不用再往下扩展。

\textbf{暴力基线。} 慢版本就是完整搜索树：每层列出选择列表，路径到达终止条件时检查答案。它先保证覆盖所有可能，再把明显无效或重复的分支剪掉。放到本题就是：先按“从一到九中选 k 个数和为 n，每个数最多用一次”完整试慢版本；样例里已经能看到“规律是路径长度、剩余和和起点共同定义状态，数字只能向后选避免重复组合”，所以优化前必须先能说清慢版本枚举了哪些候选。

\textbf{暴力过程。} 拿 k=3、n=7 手算，选择 1 后还需两个数凑 6，可以选 2 和 4；若当前和已经超过 7，就不用再往下扩展。

\textbf{重复工作。} 题面模型：从一到九中选 k 个数和为 n，每个数最多用一次。暴力版的慢点要落到本题：慢版本会围绕“从一到九中选 k 个数和为 n，每个数最多用一次”展开完整搜索树。样例规律是“规律是路径长度、剩余和和起点共同定义状态，数字只能向后选避免重复组合”，说明一层递归必须对应一个明确决策，剪枝只能删除整棵不可能成功或重复的子树。后面的优化只消掉这类重复工作，不能改变答案集合。

\textbf{优化条件。} 本题的关键观察是：按递增起点枚举，利用剩余个数和剩余和剪枝。这句话必须能翻译成可维护状态；如果题目条件改变后观察不再成立，就不能继续套原来的写法。

\textbf{相邻状态。} 规律是路径长度、剩余和和起点共同定义状态，数字只能向后选避免重复组合。把这个特例继续拆开看：一层递归对应一个明确决策，路径保存已经做出的选择，选择列表保存仍可能扩展的分支。剪枝前必须能说明被剪掉的整棵子树没有合法答案，而不是只因为它看起来不优。

\textbf{状态复用。} 把完整搜索树加上合法性检查、剩余资源剪枝和同层去重；路径状态进入和退出要成对恢复。本题落点是：规律是路径长度、剩余和和起点共同定义状态，数字只能向后选避免重复组合。手算时只改被新输入影响的那一小部分，而不是回头重算完整候选。

\textbf{指针和字段。} 状态字段要能说清：路径、起点或使用标记、剩余目标、答案集合、剪枝条件；进入递归和退出递归必须成对恢复。手算时按这个协议跟踪：拿 k=3、n=7 手算，选择 1 后还需两个数凑 6，可以选 2 和 4；若当前和已经超过 7，就不用再往下扩展。然后按协议复查：手算时给每层标出选择列表和路径。剪枝发生前要指出被剪分支为什么已经不可能得到合法答案。

\textbf{代码落点。} 剪枝条件写在扩展前；同层去重不要误写成全局去重；保存答案时复制当前路径。关键代码行必须能逐句翻译成“旧状态怎样变成新状态”，否则就只是模板。

\textbf{正确性。} 证明从搜索树覆盖开始：每个合法答案有且只有一条路径，剪枝只删掉不可能成功的分支。

\textbf{边界实验。} 先用n 太小、n 太大、k 为零、路径长度已经超过 k做反例检查。实验时覆盖n 太小、n 太大、k 为零、路径长度已经超过 k，统计搜索节点数量和剪枝次数。重复元素题要打印同一层跳过哪些值，确认没有误删合法路径。

\textbf{复杂度变化。} 暴力搜索树的规模由输出或选择空间决定；剪枝不改变最坏指数级本质，但能删除不可能成功或重复的分支，复杂度要说明输出规模。

\textbf{迁移追问。} 如果题目改成“可预估剩余最小和最大和进一步剪枝”，先判断原状态是否仍然覆盖新答案集合，再决定是微调边界、改 value 语义，还是换算法。

\noindent\textbf{LeetCode 679 24 点游戏。}

\textbf{题目说明。} LeetCode 679 24 点游戏 的题面入口要先还原成具体任务：四个数字必须全部使用，括号表达式可以看成每次选两个数合并成一个新数。

\textbf{样例入口。} 拿 [4, 1, 8, 7] 手算，先选 8 和 4 得到 4，再选 7 和 1 得到 6，最后 4*6 得到 24。若某一步选错，回溯要回到选两个数和运算符之前的状态继续试。

\textbf{暴力基线。} 慢版本就是完整搜索树：每层列出选择列表，路径到达终止条件时检查答案。它先保证覆盖所有可能，再把明显无效或重复的分支剪掉。放到本题就是：先按“四个数字必须全部使用，括号表达式可以看成每次选两个数合并成一个新数”完整试慢版本；样例里已经能看到“规律是任意括号表达式都能看成表达式树，最内层先合并两个叶子；状态是一组当前数字，每次合并两个数让集合大小减一”，所以优化前必须先能说清慢版本枚举了哪些候选。

\textbf{暴力过程。} 拿 [4, 1, 8, 7] 手算，先选 8 和 4 得到 4，再选 7 和 1 得到 6，最后 4*6 得到 24。若某一步选错，回溯要回到选两个数和运算符之前的状态继续试。

\textbf{重复工作。} 题面模型：四个数字必须全部使用，括号表达式可以看成每次选两个数合并成一个新数。暴力版的慢点要落到本题：慢版本会围绕“四个数字必须全部使用，括号表达式可以看成每次选两个数合并成一个新数”展开完整搜索树。样例规律是“规律是任意括号表达式都能看成表达式树，最内层先合并两个叶子；状态是一组当前数字，每次合并两个数让集合大小减一”，说明一层递归必须对应一个明确决策，剪枝只能删除整棵不可能成功或重复的子树。后面的优化只消掉这类重复工作，不能改变答案集合。

\textbf{优化条件。} 本题的关键观察是：从当前数字集合中枚举两个数和六种有序/无序运算，集合大小每次减少一，最后检查是否接近 24。这句话必须能翻译成可维护状态；如果题目条件改变后观察不再成立，就不能继续套原来的写法。

\textbf{相邻状态。} 规律是任意括号表达式都能看成表达式树，最内层先合并两个叶子；状态是一组当前数字，每次合并两个数让集合大小减一。把这个特例继续拆开看：一层递归对应一个明确决策，路径保存已经做出的选择，选择列表保存仍可能扩展的分支。剪枝前必须能说明被剪掉的整棵子树没有合法答案，而不是只因为它看起来不优。

\textbf{状态复用。} 把完整搜索树加上合法性检查、剩余资源剪枝和同层去重；路径状态进入和退出要成对恢复。本题落点是：规律是任意括号表达式都能看成表达式树，最内层先合并两个叶子；状态是一组当前数字，每次合并两个数让集合大小减一。手算时只改被新输入影响的那一小部分，而不是回头重算完整候选。

\textbf{指针和字段。} 状态字段要能说清：路径、起点或使用标记、剩余目标、答案集合、剪枝条件；进入递归和退出递归必须成对恢复。手算时按这个协议跟踪：拿 [4, 1, 8, 7] 手算，先选 8 和 4 得到 4，再选 7 和 1 得到 6，最后 4*6 得到 24。若某一步选错，回溯要回到选两个数和运算符之前的状态继续试。然后按协议复查：手算时给每层标出选择列表和路径。剪枝发生前要指出被剪分支为什么已经不可能得到合法答案。

\textbf{代码落点。} 剪枝条件写在扩展前；同层去重不要误写成全局去重；保存答案时复制当前路径。关键代码行必须能逐句翻译成“旧状态怎样变成新状态”，否则就只是模板。

\textbf{正确性。} 证明从搜索树覆盖开始：每个合法答案有且只有一条路径，剪枝只删掉不可能成功的分支。

\textbf{边界实验。} 先用减法和除法顺序、除零、浮点误差、重复中间状态做反例检查。实验时覆盖减法和除法顺序、除零、浮点误差、重复中间状态，统计搜索节点数量和剪枝次数。重复元素题要打印同一层跳过哪些值，确认没有误删合法路径。

\textbf{复杂度变化。} 暴力搜索树的规模由输出或选择空间决定；剪枝不改变最坏指数级本质，但能删除不可能成功或重复的分支，复杂度要说明输出规模。

\textbf{迁移追问。} 如果题目改成“可用排序后的状态去重剪枝，但核心仍是表达式树搜索”，先判断原状态是否仍然覆盖新答案集合，再决定是微调边界、改 value 语义，还是换算法。

\noindent\textbf{LeetCode 698 划分为 k 个相等的子集。}

\textbf{题目说明。} LeetCode 698 划分为 k 个相等的子集 的题面入口要先还原成具体任务：每个元素要分配到 k 个桶中，所有桶目标和相同。

\textbf{样例入口。} 拿 [4, 3, 2, 3, 5, 2, 1]、k=4 手算，总和 20，每桶目标 5；先放 5 单独成桶，再尝试 4+1、3+2。

\textbf{暴力基线。} 慢版本就是完整搜索树：每层列出选择列表，路径到达终止条件时检查答案。它先保证覆盖所有可能，再把明显无效或重复的分支剪掉。放到本题就是：先按“每个元素要分配到 k 个桶中，所有桶目标和相同”完整试慢版本；样例里已经能看到“规律是回溯状态是桶剩余容量和已用元素，排序后先放大数能更早剪枝”，所以优化前必须先能说清慢版本枚举了哪些候选。

\textbf{暴力过程。} 拿 [4, 3, 2, 3, 5, 2, 1]、k=4 手算，总和 20，每桶目标 5；先放 5 单独成桶，再尝试 4+1、3+2。

\textbf{重复工作。} 题面模型：每个元素要分配到 k 个桶中，所有桶目标和相同。暴力版的慢点要落到本题：慢版本会围绕“每个元素要分配到 k 个桶中，所有桶目标和相同”展开完整搜索树。样例规律是“规律是回溯状态是桶剩余容量和已用元素，排序后先放大数能更早剪枝”，说明一层递归必须对应一个明确决策，剪枝只能删除整棵不可能成功或重复的子树。后面的优化只消掉这类重复工作，不能改变答案集合。

\textbf{优化条件。} 本题的关键观察是：先判断总和可整除并排序降序，再用桶剩余容量回溯放置元素。这句话必须能翻译成可维护状态；如果题目条件改变后观察不再成立，就不能继续套原来的写法。

\textbf{相邻状态。} 规律是回溯状态是桶剩余容量和已用元素，排序后先放大数能更早剪枝。把这个特例继续拆开看：一层递归对应一个明确决策，路径保存已经做出的选择，选择列表保存仍可能扩展的分支。剪枝前必须能说明被剪掉的整棵子树没有合法答案，而不是只因为它看起来不优。

\textbf{状态复用。} 把完整搜索树加上合法性检查、剩余资源剪枝和同层去重；路径状态进入和退出要成对恢复。本题落点是：规律是回溯状态是桶剩余容量和已用元素，排序后先放大数能更早剪枝。手算时只改被新输入影响的那一小部分，而不是回头重算完整候选。

\textbf{指针和字段。} 状态字段要能说清：路径、起点或使用标记、剩余目标、答案集合、剪枝条件；进入递归和退出递归必须成对恢复。手算时按这个协议跟踪：拿 [4, 3, 2, 3, 5, 2, 1]、k=4 手算，总和 20，每桶目标 5；先放 5 单独成桶，再尝试 4+1、3+2。然后按协议复查：手算时给每层标出选择列表和路径。剪枝发生前要指出被剪分支为什么已经不可能得到合法答案。

\textbf{代码落点。} 剪枝条件写在扩展前；同层去重不要误写成全局去重；保存答案时复制当前路径。关键代码行必须能逐句翻译成“旧状态怎样变成新状态”，否则就只是模板。

\textbf{正确性。} 证明从搜索树覆盖开始：每个合法答案有且只有一条路径，剪枝只删掉不可能成功的分支。

\textbf{边界实验。} 先用存在元素大于目标、重复元素、空桶对称、k 为一做反例检查。实验时覆盖存在元素大于目标、重复元素、空桶对称、k 为一，统计搜索节点数量和剪枝次数。重复元素题要打印同一层跳过哪些值，确认没有误删合法路径。

\textbf{复杂度变化。} 暴力搜索树的规模由输出或选择空间决定；剪枝不改变最坏指数级本质，但能删除不可能成功或重复的分支，复杂度要说明输出规模。

\textbf{迁移追问。} 如果题目改成“状态压缩 DP 也能做，适合用掩码表示已选元素集合”，先判断原状态是否仍然覆盖新答案集合，再决定是微调边界、改 value 语义，还是换算法。

\subsection{实验复盘}

追问通常会问去重发生在同层还是路径、剪枝是否会删掉合法解、输出规模是多少。请准备画出前三层搜索树。

复盘本节时先从 LeetCode 17 电话号码的字母组合、LeetCode 22 括号生成、LeetCode 39 组合总和 开始：每道题都先手写一个能覆盖全部候选的暴力版，再指出优化结构具体保存了哪一段历史信息，最后用相邻案例比较为什么不能互换解法。

位运算和状态压缩深度题卡 按题型拆成章内大节，但每个大节都先从 LeetCode 案例切入，再进入分流、案例矩阵、案例推导和实验复盘。这样读者看到的是从具体题推到规律，而不是先读抽象方法再猜它能用在哪。

\topic{位运算与状态压缩}
这一节先看 LeetCode 29 两数相除、LeetCode 50 Pow(x, n)、LeetCode 136 只出现一次的数字 这几道 LeetCode 题，再整理同一题型的分流和推导。阅读顺序不是先背原理，而是先看案例的暴力瓶颈，再把重复出现的观察抽象成方法。

\subsection{原理和适用条件}

以 LeetCode 29 两数相除、LeetCode 50 Pow(x, n)、LeetCode 136 只出现一次的数字 为主线：LeetCode 29 两数相除：模型是“除法可以看成不断减去除数的倍增值。”，优化抓手是“把除数按二进制倍增，优先减去不超过被除数的最大倍数并累加商。”；LeetCode 50 Pow(x, n)：模型是“指数 n 可以按二进制拆分，重复乘 x 的线性过程存在大量倍增结构。”，优化抓手是“快速幂每轮根据最低位决定是否乘入答案，同时底数平方、指数右移。”；LeetCode 136 只出现一次的数字：模型是“成对元素可以用异或抵消，剩下单个元素。”，优化抓手是“异或满足交换结合，扫描顺序不影响答案。”。这些案例共同推出的规律是：先用计数或哈希保证正确，再寻找位运算的代数抵消、按位统计或函数图结构。位题的难点是证明被抵消的信息确实不影响答案。 方法名只是复盘后的标签，真正要掌握的是从题面约束、暴力失败、状态设计到边界反例的推导链。

\begin{principlebox}{图示：从 LeetCode 案例到 位运算与状态压缩推导链}
\begin{tabular}{@{}p{0.18\linewidth}p{0.74\linewidth}@{}}
暴力基线 & 枚举候选、完整模拟或逐个检查，先保证覆盖全部可能答案。\\
瓶颈定位 & 慢点在于状态复制和集合比较。位操作让包含、加入、删除、枚举子集变成常数或接近常数的机器指令。\\
结构选择 & 异或解决成对抵消，按位计数处理出现三次，掩码表示访问集合或可用集合，低位提取用于枚举候选。\\
不变量 & 掩码的每一位必须有固定含义。状态去重时，位置相同但掩码不同不能合并，因为未来能力不同。\\
代码落地 & 使用无符号整数能减少移位歧义。位数超过 31 时用 \code{long long} 或 \code{std::bitset}。表达式加括号，避免位运算优先级误读。\\
\end{tabular}
\end{principlebox}

\subsection{LeetCode 案例分流}

\begin{principlebox}{从 LeetCode 案例分流：位运算与状态压缩}
\textbf{异或抵消。}
代表案例：LeetCode 136 只出现一次的数字、LeetCode 260 只出现一次的数字 III。先看这些题的题面条件：除一个或两个数外，其他数按偶数次出现。由此推出的处理方式是：利用异或交换结合和 x 异或 x 为零。风险点：出现次数不是两次时不能直接异或。

\textbf{按位计数。}
代表案例：LeetCode 137 只出现一次的数字 II、LeetCode 剑指 Offer 56-I 数组中数字出现次数。先看这些题的题面条件：其他数字出现固定 k 次，目标数字出现不同次数。由此推出的处理方式是：逐位统计一的数量，对 k 取模还原答案。风险点：负数和符号位要使用足够宽的整数。

\textbf{掩码集合。}
代表案例：LeetCode 78 子集、LeetCode 90 子集 II、LeetCode 864 获取所有钥匙的最短路径。先看这些题的题面条件：状态是若干布尔选择或访问集合。由此推出的处理方式是：用整数位表示元素是否已选、资源是否拥有。风险点：同一位置不同掩码不能合并。

\textbf{位运算算术和低位技巧。}
代表案例：LeetCode 191 位 1 的个数、LeetCode 338 比特位计数、LeetCode 371 两整数之和。先看这些题的题面条件：题目要求不用普通算术或快速枚举子集。由此推出的处理方式是：用 lowbit、子集枚举或位加法模拟状态变化。风险点：移位优先级和有符号溢出需要谨慎。

\end{principlebox}

\subsection{案例矩阵}

本节案例覆盖 LeetCode 29 两数相除、LeetCode 50 Pow(x, n)、LeetCode 136 只出现一次的数字、LeetCode 137 只出现一次的数字 II、LeetCode 190 颠倒二进制位、LeetCode 191 位 1 的个数、LeetCode 201 数字范围按位与、LeetCode 231 2 的幂、LeetCode 260 只出现一次的数字 III、LeetCode 268 丢失的数字、LeetCode 287 寻找重复数、LeetCode 318 最大单词长度乘积、LeetCode 338 比特位计数、LeetCode 371 两整数之和、LeetCode 372 超级次方、LeetCode 393 UTF-8 编码验证、LeetCode 421 数组中两个数的最大异或值、LeetCode 1755 最接近目标值的子序列和。阅读时不要按题号背答案，而要先判断题面落在哪个分流场景：查询目标是什么，输入约束给了什么单调性或等价关系，暴力慢在哪里，优化结构保存了什么状态。

\begin{principlebox}{案例矩阵}
\textbf{LeetCode 29 两数相除。}
除法可以看成不断减去除数的倍增值。单点风险：除数为一、被除数为 int 最小值、结果溢出、符号处理。

\textbf{LeetCode 50 Pow(x, n)。}
指数 n 可以按二进制拆分，重复乘 x 的线性过程存在大量倍增结构。单点风险：n 为零、n 为负、int 最小值取反溢出、x 为零、浮点误差。

\textbf{LeetCode 136 只出现一次的数字。}
成对元素可以用异或抵消，剩下单个元素。单点风险：零、负数、唯一值在开头或结尾、所有其他值恰好两次。

\textbf{LeetCode 137 只出现一次的数字 II。}
除一个数字外，其余数字都出现三次，单纯异或不能抵消。单点风险：负数、符号位、目标为零、所有其他数恰好三次。

\textbf{LeetCode 190 颠倒二进制位。}
每一位的新位置由原位置镜像决定。单点风险：最高位、最低位、全零、全一、无符号语义。

\textbf{LeetCode 191 位 1 的个数。}
每次清掉最低位的一可以让循环次数等于一的个数。单点风险：零、只有最高位、一的个数很多、无符号输入。

\textbf{LeetCode 201 数字范围按位与。}
区间内所有数按位与后，只保留左右端公共二进制前缀。单点风险：left 等于 right、区间跨过二的幂边界、结果为零。

\textbf{LeetCode 231 2 的幂。}
二的幂在二进制中只有一个一。单点风险：n 为零、负数、int 最小值、最高位。

\textbf{LeetCode 260 只出现一次的数字 III。}
两个只出现一次的数异或后至少有一位不同。单点风险：两个答案一正一负、最低位分组、其他数恰好两次。

\textbf{LeetCode 268 丢失的数字。}
下标零到 n 与数组值相比，缺失值会在异或抵消后剩下。单点风险：缺失零、缺失 n、数组长度一、也可用求和但要防溢出。

\textbf{LeetCode 287 寻找重复数。}
数组值域使下标到值形成函数图，重复数是环入口。单点风险：重复数出现多次、不能排序、不能改数组、长度为 n 加一。

\textbf{LeetCode 318 最大单词长度乘积。}
两个单词没有公共字母时才可相乘，字母集合可压成位掩码。单点风险：重复字母、空字符串、多个单词同掩码、乘积范围。

\textbf{LeetCode 338 比特位计数。}
每个数的一的个数可由去掉最低位一后的较小数转移。单点风险：n 为零、二的幂、连续区间、表达式优先级。

\textbf{LeetCode 371 两整数之和。}
不用加号时，加法可以拆成无进位和与进位。单点风险：正负数、进位传播、有符号溢出、需要使用无符号中间值。

\textbf{LeetCode 372 超级次方。}
大指数以十进制数组给出，不能先转成普通整数。单点风险：指数为空、a 大于模数、指数含零、递归深度、模乘溢出。

\textbf{LeetCode 393 UTF-8 编码验证。}
每个字节的高位模式决定一个字符需要的后续字节数量。单点风险：非法首字节、后续字节不足、后续字节格式错误、单字节字符。

\textbf{LeetCode 421 数组中两个数的最大异或值。}
最大异或值由高位到低位尽量取一决定。单点风险：零、重复数、最高位、负数语义。

\textbf{LeetCode 1755 最接近目标值的子序列和。}
完整子序列选择是 \ensuremath{2^{n}}，拆成左右两半后各自枚举子集和。单点风险：负数、空子序列、目标为负、答案为零提前返回、两半大小。

\end{principlebox}

\subsection{共性落点}

\begin{principlebox}{本组共性落点}
\textbf{慢版本。} 暴力做法用集合、数组或递归保存状态，空间和比较成本较高。若状态本质是若干布尔选择，位掩码能把它压进整数。
\textbf{瓶颈。} 慢点在于状态复制和集合比较。位操作让包含、加入、删除、枚举子集变成常数或接近常数的机器指令。
\textbf{状态字段。} 掩码含义、当前位、目标位集合、状态转移和答案读取；负数或高位参与时要说明整数宽度。
\textbf{不变量。} 掩码的每一位必须有固定含义。状态去重时，位置相同但掩码不同不能合并，因为未来能力不同。
\textbf{实现检查。} 移位表达式加括号；使用无符号或足够宽类型；枚举子集时确认空集和全集是否包含。
\textbf{C++ 落地。} 使用无符号整数能减少移位歧义。位数超过 31 时用 \code{long long} 或 \code{std::bitset}。表达式加括号，避免位运算优先级误读。
\textbf{证明抓手。} 证明从逐位独立或子集归纳开始：被压缩到位里的信息足以决定未来转移。
\end{principlebox}

\subsection{案例推导}

\noindent\textbf{LeetCode 29 两数相除。}

\textbf{题目说明。} LeetCode 29 两数相除 的题面入口要先还原成具体任务：除法可以看成不断减去除数的倍增值。

\textbf{样例入口。} 拿 43 除以 5 手算，5 的倍增序列是 5、10、20、40，先减 40 商加 8，余 3 停止。

\textbf{暴力基线。} 慢版本先用集合、数组或布尔表保存状态，逐步执行题目动作。确认每个布尔字段的含义后，才能把它压缩成位，而不是直接写位运算公式。放到本题就是：先按“除法可以看成不断减去除数的倍增值”完整试慢版本；样例里已经能看到“规律是用二进制倍增找不超过被除数的最大除数倍数，符号和 int 最小值溢出要单独处理”，所以优化前必须先能说清慢版本枚举了哪些候选。

\textbf{暴力过程。} 拿 43 除以 5 手算，5 的倍增序列是 5、10、20、40，先减 40 商加 8，余 3 停止。

\textbf{重复工作。} 题面模型：除法可以看成不断减去除数的倍增值。暴力版的慢点要落到本题：慢版本会围绕“除法可以看成不断减去除数的倍增值”使用集合、计数或完整状态表。样例规律是“规律是用二进制倍增找不超过被除数的最大除数倍数，符号和 int 最小值溢出要单独处理”，说明哪些信息可以被逐位抵消、压缩或复用；位运算只是编码，不是证明本身。后面的优化只消掉这类重复工作，不能改变答案集合。

\textbf{优化条件。} 本题的关键观察是：把除数按二进制倍增，优先减去不超过被除数的最大倍数并累加商。这句话必须能翻译成可维护状态；如果题目条件改变后观察不再成立，就不能继续套原来的写法。

\textbf{相邻状态。} 规律是用二进制倍增找不超过被除数的最大除数倍数，符号和 int 最小值溢出要单独处理。把这个特例继续拆开看：每一位或每个掩码位都要对应题目里的一个布尔事实。位运算只是压缩表示；如果两个状态在位置相同但掩码不同，未来能力不同，就不能合并。

\textbf{状态复用。} 把集合状态压成掩码；包含、加入、删除、枚举子集都变成位操作，但每一位的题目含义不能丢。本题落点是：规律是用二进制倍增找不超过被除数的最大除数倍数，符号和 int 最小值溢出要单独处理。手算时只改被新输入影响的那一小部分，而不是回头重算完整候选。

\textbf{指针和字段。} 状态字段要能说清：掩码含义、当前位、目标位集合、状态转移和答案读取；负数或高位参与时要说明整数宽度。手算时按这个协议跟踪：拿 43 除以 5 手算，5 的倍增序列是 5、10、20、40，先减 40 商加 8，余 3 停止。然后按协议复查：手算时把整数写成二进制，并在每一位旁边标注含义。状态转移只允许改变题目动作允许改变的那些位。

\textbf{代码落点。} 移位表达式加括号；使用无符号或足够宽类型；枚举子集时确认空集和全集是否包含。关键代码行必须能逐句翻译成“旧状态怎样变成新状态”，否则就只是模板。

\textbf{正确性。} 证明从逐位独立或子集归纳开始：被压缩到位里的信息足以决定未来转移。

\textbf{边界实验。} 先用除数为一、被除数为 int 最小值、结果溢出、符号处理做反例检查。实验时覆盖除数为一、被除数为 int 最小值、结果溢出、符号处理，把数字写成二进制逐位检查。负数、最高位、空掩码和全集掩码都要单独测。

\textbf{复杂度变化。} 暴力用集合或数组保存状态可能需要更多空间和比较成本；位运算通常把每步状态更新压成常数或按位数计算，掩码 DP 仍按状态数增长。

\textbf{迁移追问。} 如果题目改成“它和快速幂类似，都利用二进制拆分减少重复加减”，先判断原状态是否仍然覆盖新答案集合，再决定是微调边界、改 value 语义，还是换算法。

\noindent\textbf{LeetCode 50 Pow(x, n)。}

\textbf{题目说明。} LeetCode 50 Pow(x, n) 的题面入口要先还原成具体任务：指数 n 可以按二进制拆分，重复乘 x 的线性过程存在大量倍增结构。

\textbf{样例入口。} 拿 x=2、n=10 手算，10 的二进制是 1010，所以答案是 \ensuremath{2^{8}} 乘 \ensuremath{2^{2}}。逐次乘 10 次太慢；每轮把底数平方、指数右移，当前位为 1 时乘入答案。

\textbf{暴力基线。} 慢版本先用集合、数组或布尔表保存状态，逐步执行题目动作。确认每个布尔字段的含义后，才能把它压缩成位，而不是直接写位运算公式。放到本题就是：先按“指数 n 可以按二进制拆分，重复乘 x 的线性过程存在大量倍增结构”完整试慢版本；样例里已经能看到“规律是指数按二进制拆分，负指数先转成倒数并用更宽整数处理”，所以优化前必须先能说清慢版本枚举了哪些候选。

\textbf{暴力过程。} 拿 x=2、n=10 手算，10 的二进制是 1010，所以答案是 \ensuremath{2^{8}} 乘 \ensuremath{2^{2}}。逐次乘 10 次太慢；每轮把底数平方、指数右移，当前位为 1 时乘入答案。

\textbf{重复工作。} 题面模型：指数 n 可以按二进制拆分，重复乘 x 的线性过程存在大量倍增结构。暴力版的慢点要落到本题：慢版本会围绕“指数 n 可以按二进制拆分，重复乘 x 的线性过程存在大量倍增结构”使用集合、计数或完整状态表。样例规律是“规律是指数按二进制拆分，负指数先转成倒数并用更宽整数处理”，说明哪些信息可以被逐位抵消、压缩或复用；位运算只是编码，不是证明本身。后面的优化只消掉这类重复工作，不能改变答案集合。

\textbf{优化条件。} 本题的关键观察是：快速幂每轮根据最低位决定是否乘入答案，同时底数平方、指数右移。这句话必须能翻译成可维护状态；如果题目条件改变后观察不再成立，就不能继续套原来的写法。

\textbf{相邻状态。} 规律是指数按二进制拆分，负指数先转成倒数并用更宽整数处理。把这个特例继续拆开看：每一位或每个掩码位都要对应题目里的一个布尔事实。位运算只是压缩表示；如果两个状态在位置相同但掩码不同，未来能力不同，就不能合并。

\textbf{状态复用。} 把集合状态压成掩码；包含、加入、删除、枚举子集都变成位操作，但每一位的题目含义不能丢。本题落点是：规律是指数按二进制拆分，负指数先转成倒数并用更宽整数处理。手算时只改被新输入影响的那一小部分，而不是回头重算完整候选。

\textbf{指针和字段。} 状态字段要能说清：掩码含义、当前位、目标位集合、状态转移和答案读取；负数或高位参与时要说明整数宽度。手算时按这个协议跟踪：拿 x=2、n=10 手算，10 的二进制是 1010，所以答案是 \ensuremath{2^{8}} 乘 \ensuremath{2^{2}}。逐次乘 10 次太慢；每轮把底数平方、指数右移，当前位为 1 时乘入答案。然后按协议复查：手算时把整数写成二进制，并在每一位旁边标注含义。状态转移只允许改变题目动作允许改变的那些位。

\textbf{代码落点。} 移位表达式加括号；使用无符号或足够宽类型；枚举子集时确认空集和全集是否包含。关键代码行必须能逐句翻译成“旧状态怎样变成新状态”，否则就只是模板。

\textbf{正确性。} 证明从逐位独立或子集归纳开始：被压缩到位里的信息足以决定未来转移。

\textbf{边界实验。} 先用n 为零、n 为负、int 最小值取反溢出、x 为零、浮点误差做反例检查。实验时覆盖n 为零、n 为负、int 最小值取反溢出、x 为零、浮点误差，把数字写成二进制逐位检查。负数、最高位、空掩码和全集掩码都要单独测。

\textbf{复杂度变化。} 暴力用集合或数组保存状态可能需要更多空间和比较成本；位运算通常把每步状态更新压成常数或按位数计算，掩码 DP 仍按状态数增长。

\textbf{迁移追问。} 如果题目改成“矩阵快速幂和模幂都是同一二进制拆分思想”，先判断原状态是否仍然覆盖新答案集合，再决定是微调边界、改 value 语义，还是换算法。

\noindent\textbf{LeetCode 136 只出现一次的数字。}

\textbf{题目说明。} LeetCode 136 只出现一次的数字 的题面入口要先还原成具体任务：成对元素可以用异或抵消，剩下单个元素。

\textbf{样例入口。} 拿 [4, 1, 2, 1, 2] 手算，1 和 1 异或抵消，2 和 2 抵消，最后剩 4。

\textbf{暴力基线。} 慢版本先用集合、数组或布尔表保存状态，逐步执行题目动作。确认每个布尔字段的含义后，才能把它压缩成位，而不是直接写位运算公式。放到本题就是：先按“成对元素可以用异或抵消，剩下单个元素”完整试慢版本；样例里已经能看到“规律是异或满足交换结合，且相同数异或后变成零，扫描顺序不影响结果；零和负数也只是按位参与同样运算”，所以优化前必须先能说清慢版本枚举了哪些候选。

\textbf{暴力过程。} 拿 [4, 1, 2, 1, 2] 手算，1 和 1 异或抵消，2 和 2 抵消，最后剩 4。

\textbf{重复工作。} 题面模型：成对元素可以用异或抵消，剩下单个元素。暴力版的慢点要落到本题：慢版本会围绕“成对元素可以用异或抵消，剩下单个元素”使用集合、计数或完整状态表。样例规律是“规律是异或满足交换结合，且相同数异或后变成零，扫描顺序不影响结果；零和负数也只是按位参与同样运算”，说明哪些信息可以被逐位抵消、压缩或复用；位运算只是编码，不是证明本身。后面的优化只消掉这类重复工作，不能改变答案集合。

\textbf{优化条件。} 本题的关键观察是：异或满足交换结合，扫描顺序不影响答案。这句话必须能翻译成可维护状态；如果题目条件改变后观察不再成立，就不能继续套原来的写法。

\textbf{相邻状态。} 规律是异或满足交换结合，且相同数异或后变成零，扫描顺序不影响结果；零和负数也只是按位参与同样运算。把这个特例继续拆开看：每一位或每个掩码位都要对应题目里的一个布尔事实。位运算只是压缩表示；如果两个状态在位置相同但掩码不同，未来能力不同，就不能合并。

\textbf{状态复用。} 把集合状态压成掩码；包含、加入、删除、枚举子集都变成位操作，但每一位的题目含义不能丢。本题落点是：规律是异或满足交换结合，且相同数异或后变成零，扫描顺序不影响结果；零和负数也只是按位参与同样运算。手算时只改被新输入影响的那一小部分，而不是回头重算完整候选。

\textbf{指针和字段。} 状态字段要能说清：掩码含义、当前位、目标位集合、状态转移和答案读取；负数或高位参与时要说明整数宽度。手算时按这个协议跟踪：拿 [4, 1, 2, 1, 2] 手算，1 和 1 异或抵消，2 和 2 抵消，最后剩 4。然后按协议复查：手算时把整数写成二进制，并在每一位旁边标注含义。状态转移只允许改变题目动作允许改变的那些位。

\textbf{代码落点。} 移位表达式加括号；使用无符号或足够宽类型；枚举子集时确认空集和全集是否包含。关键代码行必须能逐句翻译成“旧状态怎样变成新状态”，否则就只是模板。

\textbf{正确性。} 证明从逐位独立或子集归纳开始：被压缩到位里的信息足以决定未来转移。

\textbf{边界实验。} 先用零、负数、唯一值在开头或结尾、所有其他值恰好两次做反例检查。实验时覆盖零、负数、唯一值在开头或结尾、所有其他值恰好两次，把数字写成二进制逐位检查。负数、最高位、空掩码和全集掩码都要单独测。

\textbf{复杂度变化。} 暴力用集合或数组保存状态可能需要更多空间和比较成本；位运算通常把每步状态更新压成常数或按位数计算，掩码 DP 仍按状态数增长。

\textbf{迁移追问。} 如果题目改成“若其他值出现三次，要改为按位计数取模”，先判断原状态是否仍然覆盖新答案集合，再决定是微调边界、改 value 语义，还是换算法。

\noindent\textbf{LeetCode 137 只出现一次的数字 II。}

\textbf{题目说明。} LeetCode 137 只出现一次的数字 II 的题面入口要先还原成具体任务：除一个数字外，其余数字都出现三次，单纯异或不能抵消。

\textbf{样例入口。} 拿 [2, 2, 3, 2] 手算，每一位上 2 的贡献出现三次，对 3 取模后消失，只剩 3 的位。

\textbf{暴力基线。} 慢版本先用集合、数组或布尔表保存状态，逐步执行题目动作。确认每个布尔字段的含义后，才能把它压缩成位，而不是直接写位运算公式。放到本题就是：先按“除一个数字外，其余数字都出现三次，单纯异或不能抵消”完整试慢版本；样例里已经能看到“规律是逐位计数并对 3 取模，符号位也必须按固定整数宽度处理”，所以优化前必须先能说清慢版本枚举了哪些候选。

\textbf{暴力过程。} 拿 [2, 2, 3, 2] 手算，每一位上 2 的贡献出现三次，对 3 取模后消失，只剩 3 的位。

\textbf{重复工作。} 题面模型：除一个数字外，其余数字都出现三次，单纯异或不能抵消。暴力版的慢点要落到本题：慢版本会围绕“除一个数字外，其余数字都出现三次，单纯异或不能抵消”使用集合、计数或完整状态表。样例规律是“规律是逐位计数并对 3 取模，符号位也必须按固定整数宽度处理”，说明哪些信息可以被逐位抵消、压缩或复用；位运算只是编码，不是证明本身。后面的优化只消掉这类重复工作，不能改变答案集合。

\textbf{优化条件。} 本题的关键观察是：逐位统计一的数量，对三取模后还原目标数字。这句话必须能翻译成可维护状态；如果题目条件改变后观察不再成立，就不能继续套原来的写法。

\textbf{相邻状态。} 规律是逐位计数并对 3 取模，符号位也必须按固定整数宽度处理。把这个特例继续拆开看：每一位或每个掩码位都要对应题目里的一个布尔事实。位运算只是压缩表示；如果两个状态在位置相同但掩码不同，未来能力不同，就不能合并。

\textbf{状态复用。} 把集合状态压成掩码；包含、加入、删除、枚举子集都变成位操作，但每一位的题目含义不能丢。本题落点是：规律是逐位计数并对 3 取模，符号位也必须按固定整数宽度处理。手算时只改被新输入影响的那一小部分，而不是回头重算完整候选。

\textbf{指针和字段。} 状态字段要能说清：掩码含义、当前位、目标位集合、状态转移和答案读取；负数或高位参与时要说明整数宽度。手算时按这个协议跟踪：拿 [2, 2, 3, 2] 手算，每一位上 2 的贡献出现三次，对 3 取模后消失，只剩 3 的位。然后按协议复查：手算时把整数写成二进制，并在每一位旁边标注含义。状态转移只允许改变题目动作允许改变的那些位。

\textbf{代码落点。} 移位表达式加括号；使用无符号或足够宽类型；枚举子集时确认空集和全集是否包含。关键代码行必须能逐句翻译成“旧状态怎样变成新状态”，否则就只是模板。

\textbf{正确性。} 证明从逐位独立或子集归纳开始：被压缩到位里的信息足以决定未来转移。

\textbf{边界实验。} 先用负数、符号位、目标为零、所有其他数恰好三次做反例检查。实验时覆盖负数、符号位、目标为零、所有其他数恰好三次，把数字写成二进制逐位检查。负数、最高位、空掩码和全集掩码都要单独测。

\textbf{复杂度变化。} 暴力用集合或数组保存状态可能需要更多空间和比较成本；位运算通常把每步状态更新压成常数或按位数计算，掩码 DP 仍按状态数增长。

\textbf{迁移追问。} 如果题目改成“若其他数出现 k 次，思路仍是按位计数取模”，先判断原状态是否仍然覆盖新答案集合，再决定是微调边界、改 value 语义，还是换算法。

\noindent\textbf{LeetCode 190 颠倒二进制位。}

\textbf{题目说明。} LeetCode 190 颠倒二进制位 的题面入口要先还原成具体任务：每一位的新位置由原位置镜像决定。

\textbf{样例入口。} 拿低位是 101 的数手算，颠倒时每次把结果左移一位，再接上当前最低位；原数右移继续取下一位。

\textbf{暴力基线。} 慢版本先用集合、数组或布尔表保存状态，逐步执行题目动作。确认每个布尔字段的含义后，才能把它压缩成位，而不是直接写位运算公式。放到本题就是：先按“每一位的新位置由原位置镜像决定”完整试慢版本；样例里已经能看到“规律是循环 32 次处理固定宽度，不能因为高位暂时为 0 就提前结束”，所以优化前必须先能说清慢版本枚举了哪些候选。

\textbf{暴力过程。} 拿低位是 101 的数手算，颠倒时每次把结果左移一位，再接上当前最低位；原数右移继续取下一位。

\textbf{重复工作。} 题面模型：每一位的新位置由原位置镜像决定。暴力版的慢点要落到本题：慢版本会围绕“每一位的新位置由原位置镜像决定”使用集合、计数或完整状态表。样例规律是“规律是循环 32 次处理固定宽度，不能因为高位暂时为 0 就提前结束”，说明哪些信息可以被逐位抵消、压缩或复用；位运算只是编码，不是证明本身。后面的优化只消掉这类重复工作，不能改变答案集合。

\textbf{优化条件。} 本题的关键观察是：循环三十二次，把结果左移并接上当前最低位，再右移输入。这句话必须能翻译成可维护状态；如果题目条件改变后观察不再成立，就不能继续套原来的写法。

\textbf{相邻状态。} 规律是循环 32 次处理固定宽度，不能因为高位暂时为 0 就提前结束。把这个特例继续拆开看：每一位或每个掩码位都要对应题目里的一个布尔事实。位运算只是压缩表示；如果两个状态在位置相同但掩码不同，未来能力不同，就不能合并。

\textbf{状态复用。} 把集合状态压成掩码；包含、加入、删除、枚举子集都变成位操作，但每一位的题目含义不能丢。本题落点是：规律是循环 32 次处理固定宽度，不能因为高位暂时为 0 就提前结束。手算时只改被新输入影响的那一小部分，而不是回头重算完整候选。

\textbf{指针和字段。} 状态字段要能说清：掩码含义、当前位、目标位集合、状态转移和答案读取；负数或高位参与时要说明整数宽度。手算时按这个协议跟踪：拿低位是 101 的数手算，颠倒时每次把结果左移一位，再接上当前最低位；原数右移继续取下一位。然后按协议复查：手算时把整数写成二进制，并在每一位旁边标注含义。状态转移只允许改变题目动作允许改变的那些位。

\textbf{代码落点。} 移位表达式加括号；使用无符号或足够宽类型；枚举子集时确认空集和全集是否包含。关键代码行必须能逐句翻译成“旧状态怎样变成新状态”，否则就只是模板。

\textbf{正确性。} 证明从逐位独立或子集归纳开始：被压缩到位里的信息足以决定未来转移。

\textbf{边界实验。} 先用最高位、最低位、全零、全一、无符号语义做反例检查。实验时覆盖最高位、最低位、全零、全一、无符号语义，把数字写成二进制逐位检查。负数、最高位、空掩码和全集掩码都要单独测。

\textbf{复杂度变化。} 暴力用集合或数组保存状态可能需要更多空间和比较成本；位运算通常把每步状态更新压成常数或按位数计算，掩码 DP 仍按状态数增长。

\textbf{迁移追问。} 如果题目改成“若多次调用，可以按字节预处理反转表”，先判断原状态是否仍然覆盖新答案集合，再决定是微调边界、改 value 语义，还是换算法。

\noindent\textbf{LeetCode 191 位 1 的个数。}

\textbf{题目说明。} LeetCode 191 位 1 的个数 的题面入口要先还原成具体任务：每次清掉最低位的一可以让循环次数等于一的个数。

\textbf{样例入口。} 拿 11 的二进制 1011 手算，n\&(n-1) 第一次删掉最低位 1 得 1010，第二次得 1000，第三次得 0。

\textbf{暴力基线。} 慢版本先用集合、数组或布尔表保存状态，逐步执行题目动作。确认每个布尔字段的含义后，才能把它压缩成位，而不是直接写位运算公式。放到本题就是：先按“每次清掉最低位的一可以让循环次数等于一的个数”完整试慢版本；样例里已经能看到“规律是每次操作恰好删除一个 1，所以循环次数就是 1 的个数”，所以优化前必须先能说清慢版本枚举了哪些候选。

\textbf{暴力过程。} 拿 11 的二进制 1011 手算，n\&(n-1) 第一次删掉最低位 1 得 1010，第二次得 1000，第三次得 0。

\textbf{重复工作。} 题面模型：每次清掉最低位的一可以让循环次数等于一的个数。暴力版的慢点要落到本题：慢版本会围绕“每次清掉最低位的一可以让循环次数等于一的个数”使用集合、计数或完整状态表。样例规律是“规律是每次操作恰好删除一个 1，所以循环次数就是 1 的个数”，说明哪些信息可以被逐位抵消、压缩或复用；位运算只是编码，不是证明本身。后面的优化只消掉这类重复工作，不能改变答案集合。

\textbf{优化条件。} 本题的关键观察是：使用 n \& (n - 1) 删除最低位一，比逐位检查更直接。这句话必须能翻译成可维护状态；如果题目条件改变后观察不再成立，就不能继续套原来的写法。

\textbf{相邻状态。} 规律是每次操作恰好删除一个 1，所以循环次数就是 1 的个数。把这个特例继续拆开看：每一位或每个掩码位都要对应题目里的一个布尔事实。位运算只是压缩表示；如果两个状态在位置相同但掩码不同，未来能力不同，就不能合并。

\textbf{状态复用。} 把集合状态压成掩码；包含、加入、删除、枚举子集都变成位操作，但每一位的题目含义不能丢。本题落点是：规律是每次操作恰好删除一个 1，所以循环次数就是 1 的个数。手算时只改被新输入影响的那一小部分，而不是回头重算完整候选。

\textbf{指针和字段。} 状态字段要能说清：掩码含义、当前位、目标位集合、状态转移和答案读取；负数或高位参与时要说明整数宽度。手算时按这个协议跟踪：拿 11 的二进制 1011 手算，n\&(n-1) 第一次删掉最低位 1 得 1010，第二次得 1000，第三次得 0。然后按协议复查：手算时把整数写成二进制，并在每一位旁边标注含义。状态转移只允许改变题目动作允许改变的那些位。

\textbf{代码落点。} 移位表达式加括号；使用无符号或足够宽类型；枚举子集时确认空集和全集是否包含。关键代码行必须能逐句翻译成“旧状态怎样变成新状态”，否则就只是模板。

\textbf{正确性。} 证明从逐位独立或子集归纳开始：被压缩到位里的信息足以决定未来转移。

\textbf{边界实验。} 先用零、只有最高位、一的个数很多、无符号输入做反例检查。实验时覆盖零、只有最高位、一的个数很多、无符号输入，把数字写成二进制逐位检查。负数、最高位、空掩码和全集掩码都要单独测。

\textbf{复杂度变化。} 暴力用集合或数组保存状态可能需要更多空间和比较成本；位运算通常把每步状态更新压成常数或按位数计算，掩码 DP 仍按状态数增长。

\textbf{迁移追问。} 如果题目改成“比特位计数可以把这个过程 DP 化到一整段数”，先判断原状态是否仍然覆盖新答案集合，再决定是微调边界、改 value 语义，还是换算法。

\noindent\textbf{LeetCode 201 数字范围按位与。}

\textbf{题目说明。} LeetCode 201 数字范围按位与 的题面入口要先还原成具体任务：区间内所有数按位与后，只保留左右端公共二进制前缀。

\textbf{样例入口。} 拿区间 [5, 7]，二进制是 101、110、111，只有最高公共前缀 1 能留下，低位在区间内会出现 0 和 1，被按位与清零。

\textbf{暴力基线。} 慢版本先用集合、数组或布尔表保存状态，逐步执行题目动作。确认每个布尔字段的含义后，才能把它压缩成位，而不是直接写位运算公式。放到本题就是：先按“区间内所有数按位与后，只保留左右端公共二进制前缀”完整试慢版本；样例里已经能看到“规律是不断右移 left 和 right 找公共前缀，再移回原位”，所以优化前必须先能说清慢版本枚举了哪些候选。

\textbf{暴力过程。} 拿区间 [5, 7]，二进制是 101、110、111，只有最高公共前缀 1 能留下，低位在区间内会出现 0 和 1，被按位与清零。

\textbf{重复工作。} 题面模型：区间内所有数按位与后，只保留左右端公共二进制前缀。暴力版的慢点要落到本题：慢版本会围绕“区间内所有数按位与后，只保留左右端公共二进制前缀”使用集合、计数或完整状态表。样例规律是“规律是不断右移 left 和 right 找公共前缀，再移回原位”，说明哪些信息可以被逐位抵消、压缩或复用；位运算只是编码，不是证明本身。后面的优化只消掉这类重复工作，不能改变答案集合。

\textbf{优化条件。} 本题的关键观察是：不断右移 left 和 right，直到二者相等，再把公共前缀移回原位。这句话必须能翻译成可维护状态；如果题目条件改变后观察不再成立，就不能继续套原来的写法。

\textbf{相邻状态。} 规律是不断右移 left 和 right 找公共前缀，再移回原位。把这个特例继续拆开看：每一位或每个掩码位都要对应题目里的一个布尔事实。位运算只是压缩表示；如果两个状态在位置相同但掩码不同，未来能力不同，就不能合并。

\textbf{状态复用。} 把集合状态压成掩码；包含、加入、删除、枚举子集都变成位操作，但每一位的题目含义不能丢。本题落点是：规律是不断右移 left 和 right 找公共前缀，再移回原位。手算时只改被新输入影响的那一小部分，而不是回头重算完整候选。

\textbf{指针和字段。} 状态字段要能说清：掩码含义、当前位、目标位集合、状态转移和答案读取；负数或高位参与时要说明整数宽度。手算时按这个协议跟踪：拿区间 [5, 7]，二进制是 101、110、111，只有最高公共前缀 1 能留下，低位在区间内会出现 0 和 1，被按位与清零。然后按协议复查：手算时把整数写成二进制，并在每一位旁边标注含义。状态转移只允许改变题目动作允许改变的那些位。

\textbf{代码落点。} 移位表达式加括号；使用无符号或足够宽类型；枚举子集时确认空集和全集是否包含。关键代码行必须能逐句翻译成“旧状态怎样变成新状态”，否则就只是模板。

\textbf{正确性。} 证明从逐位独立或子集归纳开始：被压缩到位里的信息足以决定未来转移。

\textbf{边界实验。} 先用left 等于 right、区间跨过二的幂边界、结果为零做反例检查。实验时覆盖left 等于 right、区间跨过二的幂边界、结果为零，把数字写成二进制逐位检查。负数、最高位、空掩码和全集掩码都要单独测。

\textbf{复杂度变化。} 暴力用集合或数组保存状态可能需要更多空间和比较成本；位运算通常把每步状态更新压成常数或按位数计算，掩码 DP 仍按状态数增长。

\textbf{迁移追问。} 如果题目改成“也可以不断清除 right 的最低位一，直到 right 不大于 left”，先判断原状态是否仍然覆盖新答案集合，再决定是微调边界、改 value 语义，还是换算法。

\noindent\textbf{LeetCode 231 2 的幂。}

\textbf{题目说明。} LeetCode 231 2 的幂 的题面入口要先还原成具体任务：二的幂在二进制中只有一个一。

\textbf{样例入口。} 拿 8 的二进制 1000 手算，8\&(7)=1000\&0111=0；拿 6 的 110，6\&5=100 不为 0。

\textbf{暴力基线。} 慢版本先用集合、数组或布尔表保存状态，逐步执行题目动作。确认每个布尔字段的含义后，才能把它压缩成位，而不是直接写位运算公式。放到本题就是：先按“二的幂在二进制中只有一个一”完整试慢版本；样例里已经能看到“规律是 2 的幂只有一个 1，且 n 必须为正；0 和负数不能只靠按位与判断”，所以优化前必须先能说清慢版本枚举了哪些候选。

\textbf{暴力过程。} 拿 8 的二进制 1000 手算，8\&(7)=1000\&0111=0；拿 6 的 110，6\&5=100 不为 0。

\textbf{重复工作。} 题面模型：二的幂在二进制中只有一个一。暴力版的慢点要落到本题：慢版本会围绕“二的幂在二进制中只有一个一”使用集合、计数或完整状态表。样例规律是“规律是 2 的幂只有一个 1，且 n 必须为正；0 和负数不能只靠按位与判断”，说明哪些信息可以被逐位抵消、压缩或复用；位运算只是编码，不是证明本身。后面的优化只消掉这类重复工作，不能改变答案集合。

\textbf{优化条件。} 本题的关键观察是：正数 n 满足 n 与 n 减一按位与为零。这句话必须能翻译成可维护状态；如果题目条件改变后观察不再成立，就不能继续套原来的写法。

\textbf{相邻状态。} 规律是 2 的幂只有一个 1，且 n 必须为正；0 和负数不能只靠按位与判断。把这个特例继续拆开看：每一位或每个掩码位都要对应题目里的一个布尔事实。位运算只是压缩表示；如果两个状态在位置相同但掩码不同，未来能力不同，就不能合并。

\textbf{状态复用。} 把集合状态压成掩码；包含、加入、删除、枚举子集都变成位操作，但每一位的题目含义不能丢。本题落点是：规律是 2 的幂只有一个 1，且 n 必须为正；0 和负数不能只靠按位与判断。手算时只改被新输入影响的那一小部分，而不是回头重算完整候选。

\textbf{指针和字段。} 状态字段要能说清：掩码含义、当前位、目标位集合、状态转移和答案读取；负数或高位参与时要说明整数宽度。手算时按这个协议跟踪：拿 8 的二进制 1000 手算，8\&(7)=1000\&0111=0；拿 6 的 110，6\&5=100 不为 0。然后按协议复查：手算时把整数写成二进制，并在每一位旁边标注含义。状态转移只允许改变题目动作允许改变的那些位。

\textbf{代码落点。} 移位表达式加括号；使用无符号或足够宽类型；枚举子集时确认空集和全集是否包含。关键代码行必须能逐句翻译成“旧状态怎样变成新状态”，否则就只是模板。

\textbf{正确性。} 证明从逐位独立或子集归纳开始：被压缩到位里的信息足以决定未来转移。

\textbf{边界实验。} 先用n 为零、负数、int 最小值、最高位做反例检查。实验时覆盖n 为零、负数、int 最小值、最高位，把数字写成二进制逐位检查。负数、最高位、空掩码和全集掩码都要单独测。

\textbf{复杂度变化。} 暴力用集合或数组保存状态可能需要更多空间和比较成本；位运算通常把每步状态更新压成常数或按位数计算，掩码 DP 仍按状态数增长。

\textbf{迁移追问。} 如果题目改成“判断四的幂还要额外限制一所在的位置为偶数位”，先判断原状态是否仍然覆盖新答案集合，再决定是微调边界、改 value 语义，还是换算法。

\noindent\textbf{LeetCode 260 只出现一次的数字 III。}

\textbf{题目说明。} LeetCode 260 只出现一次的数字 III 的题面入口要先还原成具体任务：两个只出现一次的数异或后至少有一位不同。

\textbf{样例入口。} 拿 [1, 2, 1, 3, 2, 5] 手算，全部异或后只剩 3 和 5 的异或结果，取最低不同位可把 3 和 5 分到两组；每组内部重复数再异或抵消。

\textbf{暴力基线。} 慢版本先用集合、数组或布尔表保存状态，逐步执行题目动作。确认每个布尔字段的含义后，才能把它压缩成位，而不是直接写位运算公式。放到本题就是：先按“两个只出现一次的数异或后至少有一位不同”完整试慢版本；样例里已经能看到“规律是两个答案至少有一位不同，用这位分组后退化成两个只出现一次的问题”，所以优化前必须先能说清慢版本枚举了哪些候选。

\textbf{暴力过程。} 拿 [1, 2, 1, 3, 2, 5] 手算，全部异或后只剩 3 和 5 的异或结果，取最低不同位可把 3 和 5 分到两组；每组内部重复数再异或抵消。

\textbf{重复工作。} 题面模型：两个只出现一次的数异或后至少有一位不同。暴力版的慢点要落到本题：慢版本会围绕“两个只出现一次的数异或后至少有一位不同”使用集合、计数或完整状态表。样例规律是“规律是两个答案至少有一位不同，用这位分组后退化成两个只出现一次的问题”，说明哪些信息可以被逐位抵消、压缩或复用；位运算只是编码，不是证明本身。后面的优化只消掉这类重复工作，不能改变答案集合。

\textbf{优化条件。} 本题的关键观察是：用最低不同位把数组分成两组，每组内部再异或抵消。这句话必须能翻译成可维护状态；如果题目条件改变后观察不再成立，就不能继续套原来的写法。

\textbf{相邻状态。} 规律是两个答案至少有一位不同，用这位分组后退化成两个只出现一次的问题。把这个特例继续拆开看：每一位或每个掩码位都要对应题目里的一个布尔事实。位运算只是压缩表示；如果两个状态在位置相同但掩码不同，未来能力不同，就不能合并。

\textbf{状态复用。} 把集合状态压成掩码；包含、加入、删除、枚举子集都变成位操作，但每一位的题目含义不能丢。本题落点是：规律是两个答案至少有一位不同，用这位分组后退化成两个只出现一次的问题。手算时只改被新输入影响的那一小部分，而不是回头重算完整候选。

\textbf{指针和字段。} 状态字段要能说清：掩码含义、当前位、目标位集合、状态转移和答案读取；负数或高位参与时要说明整数宽度。手算时按这个协议跟踪：拿 [1, 2, 1, 3, 2, 5] 手算，全部异或后只剩 3 和 5 的异或结果，取最低不同位可把 3 和 5 分到两组；每组内部重复数再异或抵消。然后按协议复查：手算时把整数写成二进制，并在每一位旁边标注含义。状态转移只允许改变题目动作允许改变的那些位。

\textbf{代码落点。} 移位表达式加括号；使用无符号或足够宽类型；枚举子集时确认空集和全集是否包含。关键代码行必须能逐句翻译成“旧状态怎样变成新状态”，否则就只是模板。

\textbf{正确性。} 证明从逐位独立或子集归纳开始：被压缩到位里的信息足以决定未来转移。

\textbf{边界实验。} 先用两个答案一正一负、最低位分组、其他数恰好两次做反例检查。实验时覆盖两个答案一正一负、最低位分组、其他数恰好两次，把数字写成二进制逐位检查。负数、最高位、空掩码和全集掩码都要单独测。

\textbf{复杂度变化。} 暴力用集合或数组保存状态可能需要更多空间和比较成本；位运算通常把每步状态更新压成常数或按位数计算，掩码 DP 仍按状态数增长。

\textbf{迁移追问。} 如果题目改成“它是在异或抵消基础上增加分组位”，先判断原状态是否仍然覆盖新答案集合，再决定是微调边界、改 value 语义，还是换算法。

\noindent\textbf{LeetCode 268 丢失的数字。}

\textbf{题目说明。} LeetCode 268 丢失的数字 的题面入口要先还原成具体任务：下标零到 n 与数组值相比，缺失值会在异或抵消后剩下。

\textbf{样例入口。} 拿 [3, 0, 1] 手算，0、1、3 与完整下标 0、1、2、3 一起异或，重复的 0、1、3 抵消，剩 2。

\textbf{暴力基线。} 慢版本先用集合、数组或布尔表保存状态，逐步执行题目动作。确认每个布尔字段的含义后，才能把它压缩成位，而不是直接写位运算公式。放到本题就是：先按“下标零到 n 与数组值相比，缺失值会在异或抵消后剩下”完整试慢版本；样例里已经能看到“规律是缺失数可以由异或抵消得到，也可用求和但要注意溢出”，所以优化前必须先能说清慢版本枚举了哪些候选。

\textbf{暴力过程。} 拿 [3, 0, 1] 手算，0、1、3 与完整下标 0、1、2、3 一起异或，重复的 0、1、3 抵消，剩 2。

\textbf{重复工作。} 题面模型：下标零到 n 与数组值相比，缺失值会在异或抵消后剩下。暴力版的慢点要落到本题：慢版本会围绕“下标零到 n 与数组值相比，缺失值会在异或抵消后剩下”使用集合、计数或完整状态表。样例规律是“规律是缺失数可以由异或抵消得到，也可用求和但要注意溢出”，说明哪些信息可以被逐位抵消、压缩或复用；位运算只是编码，不是证明本身。后面的优化只消掉这类重复工作，不能改变答案集合。

\textbf{优化条件。} 本题的关键观察是：把所有下标和所有值一起异或，重复出现的数抵消。这句话必须能翻译成可维护状态；如果题目条件改变后观察不再成立，就不能继续套原来的写法。

\textbf{相邻状态。} 规律是缺失数可以由异或抵消得到，也可用求和但要注意溢出。把这个特例继续拆开看：每一位或每个掩码位都要对应题目里的一个布尔事实。位运算只是压缩表示；如果两个状态在位置相同但掩码不同，未来能力不同，就不能合并。

\textbf{状态复用。} 把集合状态压成掩码；包含、加入、删除、枚举子集都变成位操作，但每一位的题目含义不能丢。本题落点是：规律是缺失数可以由异或抵消得到，也可用求和但要注意溢出。手算时只改被新输入影响的那一小部分，而不是回头重算完整候选。

\textbf{指针和字段。} 状态字段要能说清：掩码含义、当前位、目标位集合、状态转移和答案读取；负数或高位参与时要说明整数宽度。手算时按这个协议跟踪：拿 [3, 0, 1] 手算，0、1、3 与完整下标 0、1、2、3 一起异或，重复的 0、1、3 抵消，剩 2。然后按协议复查：手算时把整数写成二进制，并在每一位旁边标注含义。状态转移只允许改变题目动作允许改变的那些位。

\textbf{代码落点。} 移位表达式加括号；使用无符号或足够宽类型；枚举子集时确认空集和全集是否包含。关键代码行必须能逐句翻译成“旧状态怎样变成新状态”，否则就只是模板。

\textbf{正确性。} 证明从逐位独立或子集归纳开始：被压缩到位里的信息足以决定未来转移。

\textbf{边界实验。} 先用缺失零、缺失 n、数组长度一、也可用求和但要防溢出做反例检查。实验时覆盖缺失零、缺失 n、数组长度一、也可用求和但要防溢出，把数字写成二进制逐位检查。负数、最高位、空掩码和全集掩码都要单独测。

\textbf{复杂度变化。} 暴力用集合或数组保存状态可能需要更多空间和比较成本；位运算通常把每步状态更新压成常数或按位数计算，掩码 DP 仍按状态数增长。

\textbf{迁移追问。} 如果题目改成“异或法适合展示值域完整且只有一个缺口的结构”，先判断原状态是否仍然覆盖新答案集合，再决定是微调边界、改 value 语义，还是换算法。

\noindent\textbf{LeetCode 287 寻找重复数。}

\textbf{题目说明。} LeetCode 287 寻找重复数 的题面入口要先还原成具体任务：数组值域使下标到值形成函数图，重复数是环入口。

\textbf{样例入口。} 拿 [1, 3, 4, 2, 2] 手算，把值当成 next 指针会形成 1->3->2->4->2 的环，入环点 2 就是重复数。

\textbf{暴力基线。} 慢版本先用集合、数组或布尔表保存状态，逐步执行题目动作。确认每个布尔字段的含义后，才能把它压缩成位，而不是直接写位运算公式。放到本题就是：先按“数组值域使下标到值形成函数图，重复数是环入口”完整试慢版本；样例里已经能看到“规律是长度 n+1、值域 1..n 强制形成环；不能修改数组时可用快慢指针或值域二分”，所以优化前必须先能说清慢版本枚举了哪些候选。

\textbf{暴力过程。} 拿 [1, 3, 4, 2, 2] 手算，把值当成 next 指针会形成 1->3->2->4->2 的环，入环点 2 就是重复数。

\textbf{重复工作。} 题面模型：数组值域使下标到值形成函数图，重复数是环入口。暴力版的慢点要落到本题：慢版本会围绕“数组值域使下标到值形成函数图，重复数是环入口”使用集合、计数或完整状态表。样例规律是“规律是长度 n+1、值域 1..n 强制形成环；不能修改数组时可用快慢指针或值域二分”，说明哪些信息可以被逐位抵消、压缩或复用；位运算只是编码，不是证明本身。后面的优化只消掉这类重复工作，不能改变答案集合。

\textbf{优化条件。} 本题的关键观察是：快慢指针不修改数组且常数空间；值域二分用抽屉原理。这句话必须能翻译成可维护状态；如果题目条件改变后观察不再成立，就不能继续套原来的写法。

\textbf{相邻状态。} 规律是长度 n+1、值域 1..n 强制形成环；不能修改数组时可用快慢指针或值域二分。把这个特例继续拆开看：每一位或每个掩码位都要对应题目里的一个布尔事实。位运算只是压缩表示；如果两个状态在位置相同但掩码不同，未来能力不同，就不能合并。

\textbf{状态复用。} 把集合状态压成掩码；包含、加入、删除、枚举子集都变成位操作，但每一位的题目含义不能丢。本题落点是：规律是长度 n+1、值域 1..n 强制形成环；不能修改数组时可用快慢指针或值域二分。手算时只改被新输入影响的那一小部分，而不是回头重算完整候选。

\textbf{指针和字段。} 状态字段要能说清：掩码含义、当前位、目标位集合、状态转移和答案读取；负数或高位参与时要说明整数宽度。手算时按这个协议跟踪：拿 [1, 3, 4, 2, 2] 手算，把值当成 next 指针会形成 1->3->2->4->2 的环，入环点 2 就是重复数。然后按协议复查：手算时把整数写成二进制，并在每一位旁边标注含义。状态转移只允许改变题目动作允许改变的那些位。

\textbf{代码落点。} 移位表达式加括号；使用无符号或足够宽类型；枚举子集时确认空集和全集是否包含。关键代码行必须能逐句翻译成“旧状态怎样变成新状态”，否则就只是模板。

\textbf{正确性。} 证明从逐位独立或子集归纳开始：被压缩到位里的信息足以决定未来转移。

\textbf{边界实验。} 先用重复数出现多次、不能排序、不能改数组、长度为 n 加一做反例检查。实验时覆盖重复数出现多次、不能排序、不能改数组、长度为 n 加一，把数字写成二进制逐位检查。负数、最高位、空掩码和全集掩码都要单独测。

\textbf{复杂度变化。} 暴力用集合或数组保存状态可能需要更多空间和比较成本；位运算通常把每步状态更新压成常数或按位数计算，掩码 DP 仍按状态数增长。

\textbf{迁移追问。} 如果题目改成“快慢指针要求值都能作为合法下标”，先判断原状态是否仍然覆盖新答案集合，再决定是微调边界、改 value 语义，还是换算法。

\noindent\textbf{LeetCode 318 最大单词长度乘积。}

\textbf{题目说明。} LeetCode 318 最大单词长度乘积 的题面入口要先还原成具体任务：两个单词没有公共字母时才可相乘，字母集合可压成位掩码。

\textbf{样例入口。} 拿 abcw 和 xtfn 手算，两个单词的 26 位掩码按位与为 0，说明没有公共字符，乘积 4*4 可作为候选。

\textbf{暴力基线。} 慢版本先用集合、数组或布尔表保存状态，逐步执行题目动作。确认每个布尔字段的含义后，才能把它压缩成位，而不是直接写位运算公式。放到本题就是：先按“两个单词没有公共字母时才可相乘，字母集合可压成位掩码”完整试慢版本；样例里已经能看到“规律是每个单词压成字符集合掩码，重复字母不增加位，掩码相交才排除”，所以优化前必须先能说清慢版本枚举了哪些候选。

\textbf{暴力过程。} 拿 abcw 和 xtfn 手算，两个单词的 26 位掩码按位与为 0，说明没有公共字符，乘积 4*4 可作为候选。

\textbf{重复工作。} 题面模型：两个单词没有公共字母时才可相乘，字母集合可压成位掩码。暴力版的慢点要落到本题：慢版本会围绕“两个单词没有公共字母时才可相乘，字母集合可压成位掩码”使用集合、计数或完整状态表。样例规律是“规律是每个单词压成字符集合掩码，重复字母不增加位，掩码相交才排除”，说明哪些信息可以被逐位抵消、压缩或复用；位运算只是编码，不是证明本身。后面的优化只消掉这类重复工作，不能改变答案集合。

\textbf{优化条件。} 本题的关键观察是：每个单词生成 26 位掩码，两个掩码按位与为零表示无交集。这句话必须能翻译成可维护状态；如果题目条件改变后观察不再成立，就不能继续套原来的写法。

\textbf{相邻状态。} 规律是每个单词压成字符集合掩码，重复字母不增加位，掩码相交才排除。把这个特例继续拆开看：每一位或每个掩码位都要对应题目里的一个布尔事实。位运算只是压缩表示；如果两个状态在位置相同但掩码不同，未来能力不同，就不能合并。

\textbf{状态复用。} 把集合状态压成掩码；包含、加入、删除、枚举子集都变成位操作，但每一位的题目含义不能丢。本题落点是：规律是每个单词压成字符集合掩码，重复字母不增加位，掩码相交才排除。手算时只改被新输入影响的那一小部分，而不是回头重算完整候选。

\textbf{指针和字段。} 状态字段要能说清：掩码含义、当前位、目标位集合、状态转移和答案读取；负数或高位参与时要说明整数宽度。手算时按这个协议跟踪：拿 abcw 和 xtfn 手算，两个单词的 26 位掩码按位与为 0，说明没有公共字符，乘积 4*4 可作为候选。然后按协议复查：手算时把整数写成二进制，并在每一位旁边标注含义。状态转移只允许改变题目动作允许改变的那些位。

\textbf{代码落点。} 移位表达式加括号；使用无符号或足够宽类型；枚举子集时确认空集和全集是否包含。关键代码行必须能逐句翻译成“旧状态怎样变成新状态”，否则就只是模板。

\textbf{正确性。} 证明从逐位独立或子集归纳开始：被压缩到位里的信息足以决定未来转移。

\textbf{边界实验。} 先用重复字母、空字符串、多个单词同掩码、乘积范围做反例检查。实验时覆盖重复字母、空字符串、多个单词同掩码、乘积范围，把数字写成二进制逐位检查。负数、最高位、空掩码和全集掩码都要单独测。

\textbf{复杂度变化。} 暴力用集合或数组保存状态可能需要更多空间和比较成本；位运算通常把每步状态更新压成常数或按位数计算，掩码 DP 仍按状态数增长。

\textbf{迁移追问。} 如果题目改成“若字符集大于整数位数，需要 bitset 或哈希集合”，先判断原状态是否仍然覆盖新答案集合，再决定是微调边界、改 value 语义，还是换算法。

\noindent\textbf{LeetCode 338 比特位计数。}

\textbf{题目说明。} LeetCode 338 比特位计数 的题面入口要先还原成具体任务：每个数的一的个数可由去掉最低位一后的较小数转移。

\textbf{样例入口。} 拿 n=5 手算，5 的二进制 101，比 4 多一个最低位 1，所以 bits[5]=bits[4]+1；也可用 5\&(4)=4 得到 bits[5]=bits[4]+1。

\textbf{暴力基线。} 慢版本先用集合、数组或布尔表保存状态，逐步执行题目动作。确认每个布尔字段的含义后，才能把它压缩成位，而不是直接写位运算公式。放到本题就是：先按“每个数的一的个数可由去掉最低位一后的较小数转移”完整试慢版本；样例里已经能看到“规律是每个数的答案来自去掉最低位 1 或右移一位后的更小数”，所以优化前必须先能说清慢版本枚举了哪些候选。

\textbf{暴力过程。} 拿 n=5 手算，5 的二进制 101，比 4 多一个最低位 1，所以 bits[5]=bits[4]+1；也可用 5\&(4)=4 得到 bits[5]=bits[4]+1。

\textbf{重复工作。} 题面模型：每个数的一的个数可由去掉最低位一后的较小数转移。暴力版的慢点要落到本题：慢版本会围绕“每个数的一的个数可由去掉最低位一后的较小数转移”使用集合、计数或完整状态表。样例规律是“规律是每个数的答案来自去掉最低位 1 或右移一位后的更小数”，说明哪些信息可以被逐位抵消、压缩或复用；位运算只是编码，不是证明本身。后面的优化只消掉这类重复工作，不能改变答案集合。

\textbf{优化条件。} 本题的关键观察是：dp[x] = dp[x \& (x - 1)] + 1，或 dp[x] = dp[x >> 1] + (x \& 1)。这句话必须能翻译成可维护状态；如果题目条件改变后观察不再成立，就不能继续套原来的写法。

\textbf{相邻状态。} 规律是每个数的答案来自去掉最低位 1 或右移一位后的更小数。把这个特例继续拆开看：每一位或每个掩码位都要对应题目里的一个布尔事实。位运算只是压缩表示；如果两个状态在位置相同但掩码不同，未来能力不同，就不能合并。

\textbf{状态复用。} 把集合状态压成掩码；包含、加入、删除、枚举子集都变成位操作，但每一位的题目含义不能丢。本题落点是：规律是每个数的答案来自去掉最低位 1 或右移一位后的更小数。手算时只改被新输入影响的那一小部分，而不是回头重算完整候选。

\textbf{指针和字段。} 状态字段要能说清：掩码含义、当前位、目标位集合、状态转移和答案读取；负数或高位参与时要说明整数宽度。手算时按这个协议跟踪：拿 n=5 手算，5 的二进制 101，比 4 多一个最低位 1，所以 bits[5]=bits[4]+1；也可用 5\&(4)=4 得到 bits[5]=bits[4]+1。然后按协议复查：手算时把整数写成二进制，并在每一位旁边标注含义。状态转移只允许改变题目动作允许改变的那些位。

\textbf{代码落点。} 移位表达式加括号；使用无符号或足够宽类型；枚举子集时确认空集和全集是否包含。关键代码行必须能逐句翻译成“旧状态怎样变成新状态”，否则就只是模板。

\textbf{正确性。} 证明从逐位独立或子集归纳开始：被压缩到位里的信息足以决定未来转移。

\textbf{边界实验。} 先用n 为零、二的幂、连续区间、表达式优先级做反例检查。实验时覆盖n 为零、二的幂、连续区间、表达式优先级，把数字写成二进制逐位检查。负数、最高位、空掩码和全集掩码都要单独测。

\textbf{复杂度变化。} 暴力用集合或数组保存状态可能需要更多空间和比较成本；位运算通常把每步状态更新压成常数或按位数计算，掩码 DP 仍按状态数增长。

\textbf{迁移追问。} 如果题目改成“它展示位运算和 DP 的结合，而不是逐个数循环数位”，先判断原状态是否仍然覆盖新答案集合，再决定是微调边界、改 value 语义，还是换算法。

\noindent\textbf{LeetCode 371 两整数之和。}

\textbf{题目说明。} LeetCode 371 两整数之和 的题面入口要先还原成具体任务：不用加号时，加法可以拆成无进位和与进位。

\textbf{样例入口。} 拿 5+7 手算，异或得到无进位和 2，按位与左移得到进位 10；继续迭代直到进位为 0。

\textbf{暴力基线。} 慢版本先用集合、数组或布尔表保存状态，逐步执行题目动作。确认每个布尔字段的含义后，才能把它压缩成位，而不是直接写位运算公式。放到本题就是：先按“不用加号时，加法可以拆成无进位和与进位”完整试慢版本；样例里已经能看到“规律是加法拆成无进位和与进位两部分，负数参与时要用足够明确的整数宽度处理”，所以优化前必须先能说清慢版本枚举了哪些候选。

\textbf{暴力过程。} 拿 5+7 手算，异或得到无进位和 2，按位与左移得到进位 10；继续迭代直到进位为 0。

\textbf{重复工作。} 题面模型：不用加号时，加法可以拆成无进位和与进位。暴力版的慢点要落到本题：慢版本会围绕“不用加号时，加法可以拆成无进位和与进位”使用集合、计数或完整状态表。样例规律是“规律是加法拆成无进位和与进位两部分，负数参与时要用足够明确的整数宽度处理”，说明哪些信息可以被逐位抵消、压缩或复用；位运算只是编码，不是证明本身。后面的优化只消掉这类重复工作，不能改变答案集合。

\textbf{优化条件。} 本题的关键观察是：异或得到无进位和，与运算后左移得到进位，迭代直到进位为零。这句话必须能翻译成可维护状态；如果题目条件改变后观察不再成立，就不能继续套原来的写法。

\textbf{相邻状态。} 规律是加法拆成无进位和与进位两部分，负数参与时要用足够明确的整数宽度处理。把这个特例继续拆开看：每一位或每个掩码位都要对应题目里的一个布尔事实。位运算只是压缩表示；如果两个状态在位置相同但掩码不同，未来能力不同，就不能合并。

\textbf{状态复用。} 把集合状态压成掩码；包含、加入、删除、枚举子集都变成位操作，但每一位的题目含义不能丢。本题落点是：规律是加法拆成无进位和与进位两部分，负数参与时要用足够明确的整数宽度处理。手算时只改被新输入影响的那一小部分，而不是回头重算完整候选。

\textbf{指针和字段。} 状态字段要能说清：掩码含义、当前位、目标位集合、状态转移和答案读取；负数或高位参与时要说明整数宽度。手算时按这个协议跟踪：拿 5+7 手算，异或得到无进位和 2，按位与左移得到进位 10；继续迭代直到进位为 0。然后按协议复查：手算时把整数写成二进制，并在每一位旁边标注含义。状态转移只允许改变题目动作允许改变的那些位。

\textbf{代码落点。} 移位表达式加括号；使用无符号或足够宽类型；枚举子集时确认空集和全集是否包含。关键代码行必须能逐句翻译成“旧状态怎样变成新状态”，否则就只是模板。

\textbf{正确性。} 证明从逐位独立或子集归纳开始：被压缩到位里的信息足以决定未来转移。

\textbf{边界实验。} 先用正负数、进位传播、有符号溢出、需要使用无符号中间值做反例检查。实验时覆盖正负数、进位传播、有符号溢出、需要使用无符号中间值，把数字写成二进制逐位检查。负数、最高位、空掩码和全集掩码都要单独测。

\textbf{复杂度变化。} 暴力用集合或数组保存状态可能需要更多空间和比较成本；位运算通常把每步状态更新压成常数或按位数计算，掩码 DP 仍按状态数增长。

\textbf{迁移追问。} 如果题目改成“这是位级算术模拟，和异或抵消类题目的语义不同”，先判断原状态是否仍然覆盖新答案集合，再决定是微调边界、改 value 语义，还是换算法。

\noindent\textbf{LeetCode 372 超级次方。}

\textbf{题目说明。} LeetCode 372 超级次方 的题面入口要先还原成具体任务：大指数以十进制数组给出，不能先转成普通整数。

\textbf{样例入口。} 拿 a=2、b=[1, 0] 手算，目标是 \ensuremath{2^{10}} mod 1337，不能把很长指数转成普通整数。扫描十进制指数时，新答案等于旧答案的十次方乘 a 的当前位次方。

\textbf{暴力基线。} 慢版本先用集合、数组或布尔表保存状态，逐步执行题目动作。确认每个布尔字段的含义后，才能把它压缩成位，而不是直接写位运算公式。放到本题就是：先按“大指数以十进制数组给出，不能先转成普通整数”完整试慢版本；样例里已经能看到“规律是把大指数按十进制位折叠，每步都用快速幂并取模”，所以优化前必须先能说清慢版本枚举了哪些候选。

\textbf{暴力过程。} 拿 a=2、b=[1, 0] 手算，目标是 \ensuremath{2^{10}} mod 1337，不能把很长指数转成普通整数。扫描十进制指数时，新答案等于旧答案的十次方乘 a 的当前位次方。

\textbf{重复工作。} 题面模型：大指数以十进制数组给出，不能先转成普通整数。暴力版的慢点要落到本题：慢版本会围绕“大指数以十进制数组给出，不能先转成普通整数”使用集合、计数或完整状态表。样例规律是“规律是把大指数按十进制位折叠，每步都用快速幂并取模”，说明哪些信息可以被逐位抵消、压缩或复用；位运算只是编码，不是证明本身。后面的优化只消掉这类重复工作，不能改变答案集合。

\textbf{优化条件。} 本题的关键观察是：利用 \ensuremath{(a^10)^{rest}} 和最后一位指数，把指数数组逐位折叠，并在每步取模。这句话必须能翻译成可维护状态；如果题目条件改变后观察不再成立，就不能继续套原来的写法。

\textbf{相邻状态。} 规律是把大指数按十进制位折叠，每步都用快速幂并取模。把这个特例继续拆开看：每一位或每个掩码位都要对应题目里的一个布尔事实。位运算只是压缩表示；如果两个状态在位置相同但掩码不同，未来能力不同，就不能合并。

\textbf{状态复用。} 把集合状态压成掩码；包含、加入、删除、枚举子集都变成位操作，但每一位的题目含义不能丢。本题落点是：规律是把大指数按十进制位折叠，每步都用快速幂并取模。手算时只改被新输入影响的那一小部分，而不是回头重算完整候选。

\textbf{指针和字段。} 状态字段要能说清：掩码含义、当前位、目标位集合、状态转移和答案读取；负数或高位参与时要说明整数宽度。手算时按这个协议跟踪：拿 a=2、b=[1, 0] 手算，目标是 \ensuremath{2^{10}} mod 1337，不能把很长指数转成普通整数。扫描十进制指数时，新答案等于旧答案的十次方乘 a 的当前位次方。然后按协议复查：手算时把整数写成二进制，并在每一位旁边标注含义。状态转移只允许改变题目动作允许改变的那些位。

\textbf{代码落点。} 移位表达式加括号；使用无符号或足够宽类型；枚举子集时确认空集和全集是否包含。关键代码行必须能逐句翻译成“旧状态怎样变成新状态”，否则就只是模板。

\textbf{正确性。} 证明从逐位独立或子集归纳开始：被压缩到位里的信息足以决定未来转移。

\textbf{边界实验。} 先用指数为空、a 大于模数、指数含零、递归深度、模乘溢出做反例检查。实验时覆盖指数为空、a 大于模数、指数含零、递归深度、模乘溢出，把数字写成二进制逐位检查。负数、最高位、空掩码和全集掩码都要单独测。

\textbf{复杂度变化。} 暴力用集合或数组保存状态可能需要更多空间和比较成本；位运算通常把每步状态更新压成常数或按位数计算，掩码 DP 仍按状态数增长。

\textbf{迁移追问。} 如果题目改成“若模数变化或要求精确值，取模快速幂的状态语义要重写”，先判断原状态是否仍然覆盖新答案集合，再决定是微调边界、改 value 语义，还是换算法。

\noindent\textbf{LeetCode 393 UTF-8 编码验证。}

\textbf{题目说明。} LeetCode 393 UTF-8 编码验证 的题面入口要先还原成具体任务：每个字节的高位模式决定一个字符需要的后续字节数量。

\textbf{样例入口。} 拿 [197, 130, 1] 手算，197 的高位模式表示还需要 1 个 continuation 字节，130 满足 10xxxxxx，之后 1 是单字节。

\textbf{暴力基线。} 慢版本先用集合、数组或布尔表保存状态，逐步执行题目动作。确认每个布尔字段的含义后，才能把它压缩成位，而不是直接写位运算公式。放到本题就是：先按“每个字节的高位模式决定一个字符需要的后续字节数量”完整试慢版本；样例里已经能看到“规律是状态保存剩余后续字节数量，非法首字节或后续字节不足都失败”，所以优化前必须先能说清慢版本枚举了哪些候选。

\textbf{暴力过程。} 拿 [197, 130, 1] 手算，197 的高位模式表示还需要 1 个 continuation 字节，130 满足 10xxxxxx，之后 1 是单字节。

\textbf{重复工作。} 题面模型：每个字节的高位模式决定一个字符需要的后续字节数量。暴力版的慢点要落到本题：慢版本会围绕“每个字节的高位模式决定一个字符需要的后续字节数量”使用集合、计数或完整状态表。样例规律是“规律是状态保存剩余后续字节数量，非法首字节或后续字节不足都失败”，说明哪些信息可以被逐位抵消、压缩或复用；位运算只是编码，不是证明本身。后面的优化只消掉这类重复工作，不能改变答案集合。

\textbf{优化条件。} 本题的关键观察是：用位掩码判断首字节类型，并维护还需要多少个 continuation 字节。这句话必须能翻译成可维护状态；如果题目条件改变后观察不再成立，就不能继续套原来的写法。

\textbf{相邻状态。} 规律是状态保存剩余后续字节数量，非法首字节或后续字节不足都失败。把这个特例继续拆开看：每一位或每个掩码位都要对应题目里的一个布尔事实。位运算只是压缩表示；如果两个状态在位置相同但掩码不同，未来能力不同，就不能合并。

\textbf{状态复用。} 把集合状态压成掩码；包含、加入、删除、枚举子集都变成位操作，但每一位的题目含义不能丢。本题落点是：规律是状态保存剩余后续字节数量，非法首字节或后续字节不足都失败。手算时只改被新输入影响的那一小部分，而不是回头重算完整候选。

\textbf{指针和字段。} 状态字段要能说清：掩码含义、当前位、目标位集合、状态转移和答案读取；负数或高位参与时要说明整数宽度。手算时按这个协议跟踪：拿 [197, 130, 1] 手算，197 的高位模式表示还需要 1 个 continuation 字节，130 满足 10xxxxxx，之后 1 是单字节。然后按协议复查：手算时把整数写成二进制，并在每一位旁边标注含义。状态转移只允许改变题目动作允许改变的那些位。

\textbf{代码落点。} 移位表达式加括号；使用无符号或足够宽类型；枚举子集时确认空集和全集是否包含。关键代码行必须能逐句翻译成“旧状态怎样变成新状态”，否则就只是模板。

\textbf{正确性。} 证明从逐位独立或子集归纳开始：被压缩到位里的信息足以决定未来转移。

\textbf{边界实验。} 先用非法首字节、后续字节不足、后续字节格式错误、单字节字符做反例检查。实验时覆盖非法首字节、后续字节不足、后续字节格式错误、单字节字符，把数字写成二进制逐位检查。负数、最高位、空掩码和全集掩码都要单独测。

\textbf{复杂度变化。} 暴力用集合或数组保存状态可能需要更多空间和比较成本；位运算通常把每步状态更新压成常数或按位数计算，掩码 DP 仍按状态数增长。

\textbf{迁移追问。} 如果题目改成“这是位模式解析题，重点是协议规则和状态机”，先判断原状态是否仍然覆盖新答案集合，再决定是微调边界、改 value 语义，还是换算法。

\noindent\textbf{LeetCode 421 数组中两个数的最大异或值。}

\textbf{题目说明。} LeetCode 421 数组中两个数的最大异或值 的题面入口要先还原成具体任务：最大异或值由高位到低位尽量取一决定。

\textbf{样例入口。} 拿 [3, 10, 5, 25, 2, 8] 手算，25 和 5 的异或为 28。逐位构造答案时，先假设当前位能为 1，再用哈希集合检查是否存在两个前缀配对。

\textbf{暴力基线。} 慢版本先用集合、数组或布尔表保存状态，逐步执行题目动作。确认每个布尔字段的含义后，才能把它压缩成位，而不是直接写位运算公式。放到本题就是：先按“最大异或值由高位到低位尽量取一决定”完整试慢版本；样例里已经能看到“规律是从高位到低位贪心验证前缀，不需要枚举所有数对”，所以优化前必须先能说清慢版本枚举了哪些候选。

\textbf{暴力过程。} 拿 [3, 10, 5, 25, 2, 8] 手算，25 和 5 的异或为 28。逐位构造答案时，先假设当前位能为 1，再用哈希集合检查是否存在两个前缀配对。

\textbf{重复工作。} 题面模型：最大异或值由高位到低位尽量取一决定。暴力版的慢点要落到本题：慢版本会围绕“最大异或值由高位到低位尽量取一决定”使用集合、计数或完整状态表。样例规律是“规律是从高位到低位贪心验证前缀，不需要枚举所有数对”，说明哪些信息可以被逐位抵消、压缩或复用；位运算只是编码，不是证明本身。后面的优化只消掉这类重复工作，不能改变答案集合。

\textbf{优化条件。} 本题的关键观察是：逐位尝试把当前前缀答案设为一，用哈希集合验证是否存在两个前缀配对。这句话必须能翻译成可维护状态；如果题目条件改变后观察不再成立，就不能继续套原来的写法。

\textbf{相邻状态。} 规律是从高位到低位贪心验证前缀，不需要枚举所有数对。把这个特例继续拆开看：每一位或每个掩码位都要对应题目里的一个布尔事实。位运算只是压缩表示；如果两个状态在位置相同但掩码不同，未来能力不同，就不能合并。

\textbf{状态复用。} 把集合状态压成掩码；包含、加入、删除、枚举子集都变成位操作，但每一位的题目含义不能丢。本题落点是：规律是从高位到低位贪心验证前缀，不需要枚举所有数对。手算时只改被新输入影响的那一小部分，而不是回头重算完整候选。

\textbf{指针和字段。} 状态字段要能说清：掩码含义、当前位、目标位集合、状态转移和答案读取；负数或高位参与时要说明整数宽度。手算时按这个协议跟踪：拿 [3, 10, 5, 25, 2, 8] 手算，25 和 5 的异或为 28。逐位构造答案时，先假设当前位能为 1，再用哈希集合检查是否存在两个前缀配对。然后按协议复查：手算时把整数写成二进制，并在每一位旁边标注含义。状态转移只允许改变题目动作允许改变的那些位。

\textbf{代码落点。} 移位表达式加括号；使用无符号或足够宽类型；枚举子集时确认空集和全集是否包含。关键代码行必须能逐句翻译成“旧状态怎样变成新状态”，否则就只是模板。

\textbf{正确性。} 证明从逐位独立或子集归纳开始：被压缩到位里的信息足以决定未来转移。

\textbf{边界实验。} 先用零、重复数、最高位、负数语义做反例检查。实验时覆盖零、重复数、最高位、负数语义，把数字写成二进制逐位检查。负数、最高位、空掩码和全集掩码都要单独测。

\textbf{复杂度变化。} 暴力用集合或数组保存状态可能需要更多空间和比较成本；位运算通常把每步状态更新压成常数或按位数计算，掩码 DP 仍按状态数增长。

\textbf{迁移追问。} 如果题目改成“也可以建二进制 Trie，每个数尽量走相反位”，先判断原状态是否仍然覆盖新答案集合，再决定是微调边界、改 value 语义，还是换算法。

\noindent\textbf{LeetCode 1755 最接近目标值的子序列和。}

\textbf{题目说明。} LeetCode 1755 最接近目标值的子序列和 的题面入口要先还原成具体任务：完整子序列选择是 \ensuremath{2^{n}}，拆成左右两半后各自枚举子集和。

\textbf{样例入口。} 拿 nums=[5, -7, 3, 5]、goal=6 手算，完整枚举 16 个子序列可找到 -7+3+5+5=6。折半后左半 [5, -7] 生成 0、5、-7、-2，右半 [3, 5] 生成 0、3、5、8；对每个左半和二分找最接近 goal-left 的右半和。

\textbf{暴力基线。} 慢版本先用集合、数组或布尔表保存状态，逐步执行题目动作。确认每个布尔字段的含义后，才能把它压缩成位，而不是直接写位运算公式。放到本题就是：先按“完整子序列选择是 \ensuremath{2^{n}}，拆成左右两半后各自枚举子集和”完整试慢版本；样例里已经能看到“规律是把 \ensuremath{2^{n}} 拆成两个 \ensuremath{2^{n/2}} 的子集和集合再配对”，所以优化前必须先能说清慢版本枚举了哪些候选。

\textbf{暴力过程。} 拿 nums=[5, -7, 3, 5]、goal=6 手算，完整枚举 16 个子序列可找到 -7+3+5+5=6。折半后左半 [5, -7] 生成 0、5、-7、-2，右半 [3, 5] 生成 0、3、5、8；对每个左半和二分找最接近 goal-left 的右半和。

\textbf{重复工作。} 题面模型：完整子序列选择是 \ensuremath{2^{n}}，拆成左右两半后各自枚举子集和。暴力版的慢点要落到本题：慢版本会围绕“完整子序列选择是 \ensuremath{2^{n}}，拆成左右两半后各自枚举子集和”使用集合、计数或完整状态表。样例规律是“规律是把 \ensuremath{2^{n}} 拆成两个 \ensuremath{2^{n/2}} 的子集和集合再配对”，说明哪些信息可以被逐位抵消、压缩或复用；位运算只是编码，不是证明本身。后面的优化只消掉这类重复工作，不能改变答案集合。

\textbf{优化条件。} 本题的关键观察是：分别生成两半所有子集和，排序一侧，对另一侧每个和用二分找最接近 goal 的补数。这句话必须能翻译成可维护状态；如果题目条件改变后观察不再成立，就不能继续套原来的写法。

\textbf{相邻状态。} 规律是把 \ensuremath{2^{n}} 拆成两个 \ensuremath{2^{n/2}} 的子集和集合再配对。把这个特例继续拆开看：每一位或每个掩码位都要对应题目里的一个布尔事实。位运算只是压缩表示；如果两个状态在位置相同但掩码不同，未来能力不同，就不能合并。

\textbf{状态复用。} 把集合状态压成掩码；包含、加入、删除、枚举子集都变成位操作，但每一位的题目含义不能丢。本题落点是：规律是把 \ensuremath{2^{n}} 拆成两个 \ensuremath{2^{n/2}} 的子集和集合再配对。手算时只改被新输入影响的那一小部分，而不是回头重算完整候选。

\textbf{指针和字段。} 状态字段要能说清：掩码含义、当前位、目标位集合、状态转移和答案读取；负数或高位参与时要说明整数宽度。手算时按这个协议跟踪：拿 nums=[5, -7, 3, 5]、goal=6 手算，完整枚举 16 个子序列可找到 -7+3+5+5=6。折半后左半 [5, -7] 生成 0、5、-7、-2，右半 [3, 5] 生成 0、3、5、8；对每个左半和二分找最接近 goal-left 的右半和。然后按协议复查：手算时把整数写成二进制，并在每一位旁边标注含义。状态转移只允许改变题目动作允许改变的那些位。

\textbf{代码落点。} 移位表达式加括号；使用无符号或足够宽类型；枚举子集时确认空集和全集是否包含。关键代码行必须能逐句翻译成“旧状态怎样变成新状态”，否则就只是模板。

\textbf{正确性。} 证明从逐位独立或子集归纳开始：被压缩到位里的信息足以决定未来转移。

\textbf{边界实验。} 先用负数、空子序列、目标为负、答案为零提前返回、两半大小做反例检查。实验时覆盖负数、空子序列、目标为负、答案为零提前返回、两半大小，把数字写成二进制逐位检查。负数、最高位、空掩码和全集掩码都要单独测。

\textbf{复杂度变化。} 暴力用集合或数组保存状态可能需要更多空间和比较成本；位运算通常把每步状态更新压成常数或按位数计算，掩码 DP 仍按状态数增长。

\textbf{迁移追问。} 如果题目改成“若要求方案数而非最接近差值，要改成计数配对”，先判断原状态是否仍然覆盖新答案集合，再决定是微调边界、改 value 语义，还是换算法。

\subsection{实验复盘}

追问通常会问每一位代表什么、负数和移位是否安全、状态相同但资源不同能否合并。请准备说明位运算只是编码，真正关键仍是状态语义。

复盘本节时先从 LeetCode 29 两数相除、LeetCode 50 Pow(x, n)、LeetCode 136 只出现一次的数字 开始：每道题都先手写一个能覆盖全部候选的暴力版，再指出优化结构具体保存了哪一段历史信息，最后用相邻案例比较为什么不能互换解法。

